% Problems, Chapter 7: Ordinary Differential Equations
% Arfken, Weber & Harris, Mathematical Methods for Physicists, 7th ed.
% Requires: amsmath, amssymb, tikz

\section*{Problems}

\Needspace{8\baselineskip}
\subsection*{7.2 First-Order Equations}

\Needspace{5\baselineskip}
\par\bigskip
\noindent\textbf{7.2.1}\quad From Kirchhoff's law the current $I$ in an RC
(resistance-capacitance) circuit (Fig.~7.1) obeys the equation
\[
R\frac{dI}{dt}+\frac{1}{C}I=0.
\]
\begin{enumerate}
\item[(a)] Find $I(t)$.
\item[(b)] For a capacitance of $10{,}000$ $\mu$F charged to $100$ V and
  discharging through a resistance of $1$ M$\Omega$, find the current $I$ for
  $t=0$ and for $t=100$ seconds.
\end{enumerate}
\emph{Note.}~The initial voltage is $I_{0}R$ or $Q/C$, where
$Q=\int_{0}^{\infty}I(t)\,dt$.
\solution
(a) For $R,C>0$, separate the variables and integrate from $0$ to $t$,
with $I_0=I(0)\ne0$:
\[
\frac{dI}{I}=-\frac{dt}{RC},\qquad
\ln\left|\frac{I(t)}{I_0}\right|=-\frac{t}{RC},\qquad
\boxed{I(t)=I_0e^{-t/(RC)}}.
\]
The absolute value drops because $I(t)/I_0=\pm e^{-t/(RC)}$ is continuous,
never zero, and equal to $1$ at $t=0$.
The identically zero solution is included by $I_0=0$. Differentiation
checks $RI'+I/C=-I/C+I/C=0$.

(b) Convert the given units before finding the time constant:
\[
C=10^4\times10^{-6}\ \mathrm F=10^{-2}\ \mathrm F,\quad
R=10^6\ \Omega,\quad RC=10^4\ \mathrm s.
\]
By the Note, the initial voltage is $V_0=I_0R=100$ V, so $I_0=V_0/R$;
at $t=100$ s the exponent is $t/(RC)=100/10^4=0.01$. Thus
\[
\boxed{I(0)=\frac{100}{10^6}\ \mathrm A=100\ \mu\mathrm A},\qquad
\boxed{I(100\ \mathrm s)=100e^{-0.01}\ \mu\mathrm A
=99.00498337\ldots\ \mu\mathrm A}.
\]
As a charge check,
$\int_0^\infty I(t)dt=I_0RC\bigl[-e^{-t/(RC)}\bigr]_0^\infty=I_0RC
=(10^{-4}\ \mathrm A)(10^4\ \mathrm s)=1\ \mathrm C$, which equals
$CV_0=(10^{-2}\ \mathrm F)(100\ \mathrm V)$, as the Note requires.

\Needspace{5\baselineskip}
\par\bigskip
\noindent\textbf{7.2.2}\quad The Laplace transform of Bessel's equation
($n=0$) leads to
\[
(s^{2}+1)f'(s)+sf(s)=0.
\]
Solve for $f(s)$.
\solution
For real $s$, $1+s^2$ is positive. Separate and integrate, using
$d\ln(1+s^2)/ds=2s/(1+s^2)$:
\[
\frac{f'}f=-\frac{s}{1+s^2},\qquad
\ln|f|=-\frac12\ln(1+s^2)+K.
\]
Exponentiating gives $|f|=e^K(1+s^2)^{-1/2}$; absorbing the sign into an
arbitrary real constant gives
\[
\boxed{f(s)=\frac{A}{\sqrt{1+s^2}}}.
\]
This includes $f=0$. Indeed,
$f'=-As(1+s^2)^{-3/2}$, so $(1+s^2)f'=-As(1+s^2)^{-1/2}=-sf$.
Conversely, every solution has this form, even one that vanishes
somewhere (where the logarithm fails): by the product rule,
$d\bigl[\sqrt{1+s^2}\,f\bigr]/ds=\bigl[(1+s^2)f'+sf\bigr]/\sqrt{1+s^2}=0$,
so $\sqrt{1+s^2}\,f$ is a constant $A$.
The ODE alone does not determine $A$. If the intended transform is
specifically that of $J_0$, the initial-value theorem of the Laplace
transform, $\lim_{s\to+\infty}sf(s)=J_0(0)=1$, together with
$sf(s)=As/\sqrt{1+s^2}\to A$, gives $A=1$.
For complex $s$, the same formula holds with a complex constant on a
simply connected domain avoiding $s=\pm i$, after choosing a square-root
branch. On the Laplace half-plane $\operatorname{Re}s>0$, choose the
analytic branch positive for positive real $s$; it has
$\sqrt{1+s^2}\sim s$ along the positive real axis.

\Needspace{5\baselineskip}
\par\bigskip
\noindent\textbf{7.2.3}\quad The decay of a population by catastrophic
two-body collisions is described by
\[
\frac{dN}{dt}=-kN^{2}.
\]
This is a first-order, nonlinear differential equation. Derive the solution
\[
N(t)=N_{0}\left(1+\frac{t}{\tau_{0}}\right)^{-1},
\]
where $\tau_{0}=(kN_{0})^{-1}$, \erratum{with $k,N_0>0$.
The backward extrapolation has a pole as $t\downarrow-\tau_0$;
$N(-\tau_0)$ itself is undefined.}
\solution
For $N_0>0$ and $k>0$, integrate from the initial state:
\[
\int_{N_0}^{N(t)}\frac{dN}{N^2}=-k\int_0^t ds,
\qquad -\frac1{N(t)}+\frac1{N_0}=-kt.
\]
Hence $1/N(t)=1/N_0+kt=(1+kN_0t)/N_0$, and taking the reciprocal gives
\[
\boxed{N(t)=\frac{N_0}{1+kN_0t}
=N_0\left(1+\frac t{\tau_0}\right)^{-1}},\qquad
\tau_0=\frac1{kN_0}.
\]
Direct differentiation yields $N'=-kN_0^2/(1+kN_0t)^2=-kN^2$.
The initial-value solution with positive population has mathematical
interval $t> -\tau_0$ and a pole at its left endpoint; physical decay
uses $t\ge0$. The textbook's closing remark that this implies an
infinite population at $t=-\tau_0$, qualified by the red note in the
statement, should be read accordingly: the pole is a backward
extrapolation of the model, not an attained infinite population at a
finite time within its solution interval.
If $N_0=0$, the separate solution is $N\equiv0$; if $k=0$, $N\equiv N_0$.

\Needspace{5\baselineskip}
\par\bigskip
\noindent\textbf{7.2.4}\quad The rate of a particular chemical reaction
$\mathrm{A}+\mathrm{B}\to\mathrm{C}$ is proportional to the concentrations of
the reactants A and B:
\[
\frac{dC(t)}{dt}=\alpha[A(0)-C(t)][B(0)-C(t)].
\]
\begin{enumerate}
\item[(a)] Find $C(t)$ for $A(0)\neq B(0)$.
\item[(b)] Find $C(t)$ for $A(0)=B(0)$.
\end{enumerate}
The initial condition is that $C(0)=0$.
\solution
Let $A_0=A(0)$ and $B_0=B(0)$, and assume nonnegative initial
concentrations and $\alpha\ge0$.

(a) If $A_0\ne B_0$ and both are positive, partial fractions give
\[
\frac1{(A_0-C)(B_0-C)}
=\frac1{B_0-A_0}\left(\frac1{A_0-C}-\frac1{B_0-C}\right),
\]
as recombining the right side over the common denominator shows.
Separate the variables, $dC/[(A_0-C)(B_0-C)]=\alpha\,dt$, and integrate
from $0$ to $t$. The constant $C=\min(A_0,B_0)$ is an equilibrium
solution. The right side $\alpha(A_0-C)(B_0-C)$ is a polynomial in $C$,
so by the uniqueness theorem for $C'=f(C)$ with continuously
differentiable $f$ two solutions cannot take the same value at the same
time; hence the solution starting at $C(0)=0$ never reaches the
equilibrium. Below it both factors are positive, so $C'\ge0$ and $C$
cannot fall below $C(0)=0$. Thus $0\le C<\min(A_0,B_0)$, where
\[
\int_0^C\frac{du}{A_0-u}=\ln\frac{A_0}{A_0-C},\qquad
\int_0^C\frac{du}{B_0-u}=\ln\frac{B_0}{B_0-C}.
\]
The difference of these logarithms, divided by $B_0-A_0$, yields
\[
\frac1{B_0-A_0}\ln\left[\frac{B_0-C}{A_0-C}\frac{A_0}{B_0}\right]
=\alpha t.
\]
Set $E=e^{\alpha(B_0-A_0)t}$. Multiplying by $B_0-A_0$, exponentiating,
clearing the denominators, and collecting the terms in $C$ give
\[
A_0(B_0-C)=B_0E(A_0-C),\qquad
C(B_0E-A_0)=A_0B_0(E-1).
\]
Thus $C=A_0B_0(E-1)/(B_0E-A_0)$; dividing numerator and denominator by
$E$ gives
\[
\boxed{C(t)=\frac{A_0B_0(1-e^{-\alpha(B_0-A_0)t})}
{B_0-A_0e^{-\alpha(B_0-A_0)t}}}.
\]
It starts at zero and approaches $\min(A_0,B_0)$ for $\alpha>0$: if
$B_0>A_0$ the exponential tends to $0$ and $C\to A_0B_0/B_0=A_0$; if
$B_0<A_0$ it grows without bound, and $C\to A_0B_0/A_0=B_0$.
To check the result, put $\lambda=\alpha(B_0-A_0)$ and
$D=B_0-A_0e^{-\lambda t}$, so $C=A_0B_0(1-e^{-\lambda t})/D$. Then
$A_0-C=A_0(B_0-A_0)e^{-\lambda t}/D$ and $B_0-C=B_0(B_0-A_0)/D$, so
$\alpha(A_0-C)(B_0-C)=\alpha A_0B_0(B_0-A_0)^2e^{-\lambda t}/D^2$.
Since $D'=\lambda A_0e^{-\lambda t}$, the quotient rule gives
\[
C'=\frac{A_0B_0\lambda e^{-\lambda t}\bigl[D-A_0(1-e^{-\lambda t})\bigr]}{D^2}
=\frac{A_0B_0\lambda e^{-\lambda t}(B_0-A_0)}{D^2},
\]
the same expression because $\lambda=\alpha(B_0-A_0)$.
If either concentration is zero, $C(t)=0$ directly; the formula also
reduces to that value when $A_0\ne B_0$.

(b) With $A_0=B_0=a>0$, separate instead as
\[
\int_0^{C(t)}\frac{du}{(a-u)^2}=\alpha t,
\qquad \frac1{a-C}-\frac1a=\alpha t.
\]
Hence $1/(a-C)=(1+\alpha at)/a$, i.e., $a-C=a/(1+\alpha at)$, and
\[
\boxed{C(t)=a-\frac a{1+\alpha at}=\frac{\alpha a^2t}{1+\alpha at}}.
\]
The formula includes $a=0$ and $\alpha=0$. It satisfies the equation:
$C'=\alpha a^2/(1+\alpha at)^2=\alpha(a-C)^2$.

\Needspace{5\baselineskip}
\par\bigskip
\noindent\textbf{7.2.5}\quad A boat, coasting through the water, experiences a
resisting force proportional to $v^{n}$, $v$ being the boat's instantaneous
velocity. Newton's second law leads to
\[
m\frac{dv}{dt}=-kv^{n}.
\]
With $v(t=0)=v_{0}$, $x(t=0)=0$, integrate to find $v$ as a function of time
and $v$ as a function of distance.
\erratum{Take $m,k,v_0>0$, $n\in\mathbb R$, and use the real positive-speed
branch until it reaches zero, if it does.}
\solution
Take $m,k,v_0>0$ and interpret $v$ as speed while the boat moves
forward. Put $\gamma=k/m$. For $n\ne1$, separate and integrate from $0$
to $t$:
\[
v^{-n}dv=-\gamma dt,\qquad
\frac{v^{1-n}-v_0^{1-n}}{1-n}=-\gamma t,
\]
so $v^{1-n}=v_0^{1-n}-(1-n)\gamma t$, and raising to the power $1/(1-n)$,
\[
\boxed{v(t)=\left[v_0^{1-n}+(n-1)\gamma t\right]^{1/(1-n)}}
\quad(n\ne1).
\]
The real positive branch is used while the bracket is positive. For
$n=1$, the integration is logarithmic:
\[
\ln(v/v_0)=-\gamma t,\qquad \boxed{v(t)=v_0e^{-\gamma t}}.
\]
To eliminate time use $dv/dt=(dv/dx)(dx/dt)=v\,dv/dx$ while $v>0$.
Then $mv\,dv/dx=-kv^n$, i.e., $v^{1-n}dv=-\gamma dx$; integrating from
$x=0$, for $n\ne2$,
\[
\frac{v^{2-n}-v_0^{2-n}}{2-n}=-\gamma x,
\qquad
\boxed{v(x)=\left[v_0^{2-n}+(n-2)\gamma x\right]^{1/(2-n)}}.
\]
For the omitted case $n=2$, $dv/v=-\gamma dx$, giving
\[
\boxed{v(x)=v_0e^{-\gamma x}}.
\]
The general formulas also cover the value excluded in the other
variable: $n=1$ in $v(x)$ gives $v(x)=v_0-\gamma x$, and $n=2$ in $v(t)$
gives $v(t)=(v_0^{-1}+\gamma t)^{-1}=v_0/(1+\gamma v_0t)$.
Differentiating the two integrated identities,
$d[v^{1-n}/(1-n)]/dt=v^{-n}v'=-\gamma$ and
$d[v^{2-n}/(2-n)]/dx=v^{1-n}\,dv/dx=-\gamma$, checks both equations
for every $n\ne1$ and every $n\ne2$, respectively; the exponentials
satisfy $v'=-\gamma v$ ($n=1$) and $v\,dv/dx=-\gamma v^2$ ($n=2$)
directly.

For $n<1$ the bracket in $v(t)$ decreases to zero, so the positive-speed
branch reaches zero at $t_*=v_0^{1-n}/[(1-n)\gamma]$; likewise, for $n<2$
the bracket in $v(x)$ vanishes at the limiting travel distance
$x_*=v_0^{2-n}/[(2-n)\gamma]$. For $n>1$ the bracket in $v(t)$ grows
without bound, so $v\to0$ only as $t\to\infty$ (as for the exponential
at $n=1$): stopping takes infinite time. In the same way, for $n\ge2$ the
speed $v(x)$ is positive at every $x$, so the traveled distance is
unbounded: were $x(t)\le L$ for all $t$, then, since
$dv/dx=-\gamma v^{n-1}<0$, the speed would stay at least $v(L)>0$, so
$x(t)\ge v(L)\,t$ would exceed $L$ for $t>L/v(L)$, a contradiction.
\erratum{The positive-speed drag model does not authorize continuing a
fractional or negative power through $v=0$. At or after stopping an
appropriate force law must be specified; the formulas above state their
actual domain.}

\Needspace{5\baselineskip}
\par\bigskip
\noindent\textbf{7.2.6}\quad In the first-order differential equation
$dy/dx=f(x,y)$, the function $f(x,y)$ is a function of the ratio $y/x$:
\[
\frac{dy}{dx}=g(y/x).
\]
Show that the substitution of $u=y/x$ leads to a separable equation in $u$ and
$x$.
\solution
Work on an interval with $x\ne0$. Since $y=xu$,
\[
\frac{dy}{dx}=u+x\frac{du}{dx}=g(u),\qquad
x\frac{du}{dx}=g(u)-u.
\]
Thus, whenever $g(u)-u\ne0$,
\[
\boxed{\frac{du}{g(u)-u}=\frac{dx}{x}},\qquad
\int^u\frac{ds}{g(s)-s}=\ln|x|+C.
\]
Division by $g(u)-u$ omits its zeros. Each constant $u_*$ satisfying
$g(u_*)=u_*$ supplies the additional straight-line solution $y=u_*x$,
which must be checked before declaring the solution set complete.

\Needspace{5\baselineskip}
\par\bigskip
\noindent\textbf{7.2.7}\quad The differential equation
\[
P(x,y)\,dx+Q(x,y)\,dy=0
\]
is exact. \erratum{Assume $P,Q$ have continuous first partial derivatives
and that the rectangular integration path lies inside the domain.} Show that its solution is of the form
\[
\varphi(x,y)=\erratum{\int_{x_{0}}^{x}P(\xi,y)\,d\xi+\int_{y_{0}}^{y}Q(x_{0},\eta)\,d\eta}
=\text{constant}.
\]
\solution
Use distinct integration variables, so the proposed potential is
\[
\varphi(x,y)=\int_{x_0}^xP(\xi,y)d\xi
             +\int_{y_0}^yQ(x_0,\eta)d\eta.
\]
This is the line integral along the path first from $(x_0,y_0)$ to
$(x_0,y)$, then to $(x,y)$. Exactness means $P\,dx+Q\,dy=d\Phi$
for a potential $\Phi$, i.e., $\Phi_x=P$ and $\Phi_y=Q$. Hence, by the
fundamental theorem of calculus,
\[
\int_{y_0}^yQ(x_0,\eta)d\eta=\Phi(x_0,y)-\Phi(x_0,y_0),\qquad
\int_{x_0}^xP(\xi,y)d\xi=\Phi(x,y)-\Phi(x_0,y),
\]
and the sum is $\varphi(x,y)=\Phi(x,y)-\Phi(x_0,y_0)$. It follows that
$d\varphi=d\Phi=P\,dx+Q\,dy$ and the solution curves satisfy
\[
\boxed{\varphi(x,y)=C}.
\]
More explicitly, differentiation (under the integral sign in the first
term, then the exactness condition $P_y=Q_x$ and the fundamental theorem
of calculus in $\xi$; see 7.2.8) gives
\[
\varphi_x=P(x,y),\quad
\varphi_y=\int_{x_0}^xP_y(\xi,y)d\xi+Q(x_0,y)
=Q(x,y)-Q(x_0,y)+Q(x_0,y)=Q(x,y).
\]
Along any solution curve $(x(t),y(t))$, the chain rule gives
$d\varphi/dt=P\dot x+Q\dot y=0$. Conversely this same calculation
shows that a differentiable level curve solves the differential form.
The displayed rectangular path is understood to lie inside the domain;
otherwise one uses an admissible path for the same exact differential.

\Needspace{5\baselineskip}
\par\bigskip
\noindent\textbf{7.2.8}\quad The differential equation
\[
P(x,y)\,dx+Q(x,y)\,dy=0
\]
is exact. \erratum{Assume $P,Q$ have continuous first partial derivatives
and that the rectangular integration path lies inside the domain.} If
\[
\varphi(x,y)=\erratum{\int_{x_{0}}^{x}P(\xi,y)\,d\xi+\int_{y_{0}}^{y}Q(x_{0},\eta)\,d\eta},
\]
show that
\[
\frac{\partial\varphi}{\partial x}=P(x,y),\quad
\frac{\partial\varphi}{\partial y}=Q(x,y).
\]
Hence, $\varphi(x,y)=\text{constant}$ is a solution of the original
differential equation.
\solution
Write the formula with dummy variables:
\[
\varphi(x,y)=\int_{x_0}^xP(\xi,y)d\xi
             +\int_{y_0}^yQ(x_0,\eta)d\eta.
\]
The fundamental theorem of calculus, with $y$ held fixed, gives
$\boxed{\varphi_x=P(x,y)}$; the second term does not depend on $x$.
Differentiating with respect to $y$ requires both terms: the first is
differentiated under the integral sign (allowed because $P_y$ is
continuous), the second by the fundamental theorem of calculus in $\eta$:
\begin{align*}
\varphi_y
&=\int_{x_0}^xP_y(\xi,y)d\xi+Q(x_0,y)\\
&=\int_{x_0}^xQ_x(\xi,y)d\xi+Q(x_0,y)\\
&=Q(x,y)-Q(x_0,y)+Q(x_0,y)=\boxed{Q(x,y)}.
\end{align*}
The second line uses the exactness condition $P_y=Q_x$; the third
integrates $Q_x$ in $\xi$ by the fundamental theorem of calculus.
Therefore $d\varphi=\varphi_xdx+\varphi_ydy=Pdx+Qdy$.
A level curve $\varphi=C$ has $d\varphi=0$ and satisfies the original
equation. As in 7.2.7, the chosen paths must lie within the domain.

\Needspace{5\baselineskip}
\par\bigskip
\noindent\textbf{7.2.9}\quad Prove that Eq.~(7.12) is exact in the sense of
Eq.~(7.9), provided that $\alpha(x)$ satisfies Eq.~(7.14).
\solution
The referenced equations are
$\alpha y'+\alpha py=\alpha q$ and $\alpha'=\alpha p$.
In differential form (multiply the first by $dx$ and move $\alpha q\,dx$
to the left) it is
\[
M(x,y)dx+N(x,y)dy=0,\qquad
M=\alpha(x)[p(x)y-q(x)],\quad N=\alpha(x).
\]
The exactness test is
\[
M_y=\alpha p,\qquad N_x=\alpha',
\]
so $M_y=N_x$ follows precisely from Eq.~(7.14). A direct potential is
\[
\boxed{\Phi(x,y)=\alpha(x)y-\int^x\alpha(s)q(s)ds},
\]
because $\Phi_y=\alpha=N$ and, by Eq.~(7.14),
$\Phi_x=\alpha' y-\alpha q=\alpha(py-q)=M$.
For equivalence with the original ODE choose a nonzero integrating factor,
$\alpha=e^{\int p\,dx}$; the identically zero factor would give only
$0=0$ and discard the equation.

\Needspace{5\baselineskip}
\par\bigskip
\noindent\textbf{7.2.10}\quad A certain differential equation has the form
\[
f(x)\,dx+g(x)h(y)\,dy=0,
\]
with none of the functions $f(x)$, $g(x)$, $h(y)$ identically zero. Show that
a necessary and sufficient condition for this equation to be exact is that
$g(x)=\text{constant}$.
\erratum{Work on a product of connected intervals, with $C^1$ coefficients.}
\solution
On a product of connected intervals, with the stated functions continuously
differentiable where needed, exactness requires
\[
\frac{\partial f(x)}{\partial y}=0
=\frac{\partial[g(x)h(y)]}{\partial x}=g'(x)h(y).
\]
Since $h$ is not identically zero, choose one $y_*$ with $h(y_*)\ne0$.
The equality at $(x,y_*)$ forces $g'(x)=0$ for every $x$, so $g$ is a
constant. Conversely, if $g=g_0$, then
\[
f(x)dx+g_0h(y)dy=d\left[\int^xf(s)ds+g_0\int^yh(t)dt\right]
\]
is exact. Hence $\boxed{g(x)=g_0}$ is necessary and sufficient.
\erratum{The connected product-domain assumption specifies the usual
meaning of this exercise. Without it, $h$ could vanish throughout an
entire component, imposing no condition on $g$ there, or different
$x$-components could carry different constants.}

\Needspace{5\baselineskip}
\par\bigskip
\noindent\textbf{7.2.11}\quad Show that
\[
y(x)=\exp\left[-\int^{x}p(t)\,dt\right]
\left\{\int^{x}\left[\exp\left(\int^{s}p(t)\,dt\right)\right]q(s)\,ds
+C\right\}
\]
is a solution of
\[
\frac{dy}{dx}+p(x)y(x)=q(x)
\]
by differentiating the expression for $y(x)$ and substituting into the
differential equation.
\solution
Choose a fixed base point and abbreviate
\[
P_0(x)=\int_{x_0}^xp(t)dt,\qquad
J(x)=\int_{x_0}^xe^{P_0(s)}q(s)ds+C.
\]
Then the proposed solution is $y=e^{-P_0}J$, with
$P_0'=p$ and $J'=e^{P_0}q$. The product and chain rules give
\[
y'=-p e^{-P_0}J+e^{-P_0}e^{P_0}q=-py+q.
\]
Therefore $\boxed{y'+py=q}$. Changing the fixed lower limits merely
changes the arbitrary constant or rescales the integrating factor and
does not change the resulting family of solutions.

\Needspace{5\baselineskip}
\par\bigskip
\noindent\textbf{7.2.12}\quad The motion of a body falling in a resisting
medium may be described by
\[
m\frac{dv}{dt}=mg-bv
\]
when the retarding force is proportional to the velocity, $v$. Find the
velocity. Evaluate the constant of integration by demanding that $v(0)=0$.
\solution
For $b>0$, divide by $m$ and multiply by the integrating factor
$e^{bt/m}$; by the product rule,
$(e^{bt/m}v)'=e^{bt/m}(v'+bv/m)$, so
\[
v'+\frac bm v=g,\qquad
\frac d{dt}(e^{bt/m}v)=ge^{bt/m}.
\]
Integrating from $0$ to $t$, using $v(0)=0$ and
$\int_0^tge^{bs/m}ds=(mg/b)(e^{bt/m}-1)$, gives
\[
e^{bt/m}v(t)=\frac{mg}{b}(e^{bt/m}-1),\qquad
\boxed{v(t)=\frac{mg}{b}(1-e^{-bt/m})}.
\]
The check is $mv'=mg e^{-bt/m}=mg-bv$.
The terminal velocity is $mg/b$; as $b\to0$, the expansion
$1-e^{-bt/m}=bt/m+O(b^2)$ gives the drag-free solution $v=gt$.

\Needspace{5\baselineskip}
\par\bigskip
\noindent\textbf{7.2.13}\quad Radioactive nuclei decay according to the law
\[
\frac{dN}{dt}=-\lambda N,
\]
$N$ being the concentration of a given nuclide and $\lambda$, the particular
decay constant. In a radioactive series of two different nuclides, with
concentrations $N_{1}(t)$ and $N_{2}(t)$, we have
\[
\frac{dN_{1}}{dt}=-\lambda_{1}N_{1},\quad
\frac{dN_{2}}{dt}=\lambda_{1}N_{1}-\lambda_{2}N_{2}.
\]
Find $N_{2}(t)$ for the conditions $N_{1}(0)=N_{0}$ and $N_{2}(0)=0$.
\solution
The first equation integrates to $N_1=N_0e^{-\lambda_1t}$. The second
becomes a linear equation; multiplying it by the integrating factor
$e^{\lambda_2t}$ and using
$(e^{\lambda_2t}N_2)'=e^{\lambda_2t}(N_2'+\lambda_2N_2)$,
\[
N_2'+\lambda_2N_2=\lambda_1N_0e^{-\lambda_1t},\qquad
\frac d{dt}(e^{\lambda_2t}N_2)
=\lambda_1N_0e^{(\lambda_2-\lambda_1)t}.
\]
Integrate from $0$ to $t$, with $N_2(0)=0$. If
$\lambda_1\ne\lambda_2$,
\[
e^{\lambda_2t}N_2(t)=\lambda_1N_0\int_0^te^{(\lambda_2-\lambda_1)s}ds
=\frac{\lambda_1N_0}{\lambda_2-\lambda_1}
\bigl(e^{(\lambda_2-\lambda_1)t}-1\bigr),
\]
and multiplication by $e^{-\lambda_2t}$ gives
\[
\boxed{N_2(t)=\frac{\lambda_1N_0}{\lambda_2-\lambda_1}
(e^{-\lambda_1t}-e^{-\lambda_2t})}.
\]
\erratum{The coincident-rate case must be handled separately rather than
dividing by $\lambda_2-\lambda_1$.} If both rates equal $\lambda$,
\[
e^{\lambda t}N_2=\lambda N_0\int_0^t ds=\lambda N_0t,
\qquad \boxed{N_2=\lambda N_0t e^{-\lambda t}}.
\]
This is also the limit of the preceding expression: with
$\lambda_1=\lambda$ and $\lambda_2\to\lambda$,
$(e^{-\lambda t}-e^{-\lambda_2t})/(\lambda_2-\lambda)$ tends to the
derivative of $-e^{-\lambda_2t}$ with respect to $\lambda_2$ at
$\lambda_2=\lambda$, which is $te^{-\lambda t}$. Both forms satisfy
the initial condition and the inhomogeneous equation by differentiation.

\Needspace{5\baselineskip}
\par\bigskip
\noindent\textbf{7.2.14}\quad The rate of evaporation from a particular
spherical drop of liquid (constant density) is proportional to its surface
area. Assuming this to be the sole mechanism of mass loss, find the radius of
the drop as a function of time.
\solution
Let $\rho$ be the fixed mass density and let $j>0$ be the mass lost per
unit area per unit time. Then
\[
M=\frac{4\pi}3\rho r^3,\qquad
\frac{dM}{dt}=4\pi\rho r^2r'=-j(4\pi r^2).
\]
While $r>0$, cancellation gives $r'=-j/\rho$. With $r(0)=r_0$,
\[
\boxed{r(t)=r_0-\frac j\rho t,\qquad
0\le t\le t_*:=\frac{\rho r_0}{j}.}
\]
At $t_*$ the drop is exhausted. The model's physical continuation has
zero radius and mass, not a negative radius. If the proportionality
constant is defined for volume loss instead of mass loss, put
$\kappa=j/\rho$ and write $r=r_0-\kappa t$.

\Needspace{5\baselineskip}
\par\bigskip
\noindent\textbf{7.2.15}\quad In the linear homogeneous differential equation
\[
\frac{dv}{dt}=-av
\]
the variables are separable. When the variables are separated, the equation is
exact. Solve this differential equation subject to $v(0)=v_{0}$ by the
following three methods:
\begin{enumerate}
\item[(a)] Separating variables and integrating.
\item[(b)] Treating the separated variable equation as exact.
\item[(c)] Using the result for a linear homogeneous differential equation.
\end{enumerate}
\par\nobreak\hfill\emph{ANS.}\enspace $v(t)=v_{0}e^{-at}$.
\solution
(a) For $v_0\ne0$, separation gives
\[
\int_{v_0}^{v(t)}\frac{du}{u}=-a\int_0^t ds,
\quad \ln|v(t)/v_0|=-at,
\quad v(t)=v_0e^{-at}.
\]
The sign stays that of $v_0$; $v\equiv0$ is separately included when
$v_0=0$.

(b) The separated form $a\,dt+dv/v=0$, i.e., $P\,dt+Q\,dv=0$ with
$P=a$ and $Q=1/v$, is exact on either half-plane $v>0$ or $v<0$, since
both cross partial derivatives vanish: $\partial P/\partial v=0
=\partial Q/\partial t$. Its potential is $\Phi(t,v)=at+\ln|v|$, because
$\Phi_t=a$ and $\Phi_v=1/v$. The level condition
$\Phi(t,v)=\Phi(0,v_0)=\ln|v_0|$ gives $\ln|v/v_0|=-at$, so
$|v|=|v_0|e^{-at}$; since $v$ stays in the half-plane of $v_0$, this is
the same exponential $v=v_0e^{-at}$.
The zero solution was excluded only by division by $v$ and must be restored.

(c) In the linear first-order formula of Exercise 7.2.11 take $p=a$,
$q=0$:
\[
v=e^{-\int a\,dt}(C+0)=Ce^{-at}.
\]
The initial condition gives $C=v_0$. Thus all three methods give
$\boxed{v=v_0e^{-at}}$, and $v'=-av$ checks it directly.

\Needspace{5\baselineskip}
\par\bigskip
\noindent\textbf{7.2.16}\quad
\begin{enumerate}
\item[(a)] Solve Example 7.2.1, assuming that the parachute opens when the
  parachutist's velocity has reached $v_{i}=60$ mi/h (regard this time as
  $t=0$). Find $v(t)$.
\item[(b)] For a skydiver in free fall use the friction coefficient
  $b=0.25$ kg/m and mass $m=70$ kg. What is the limiting velocity in this
  case?
\end{enumerate}
\solution
In Example 7.2.1 the downward speed obeys
$mv'=mg-bv^2$. Put $v_T=\sqrt{mg/b}$ and
$T=\sqrt{m/(bg)}$, so $v_T/T=g$.

(a) Since $mg-bv^2=b(v_T^2-v^2)$, separation gives
$dv/(v_T^2-v^2)=(b/m)\,dt$. By partial fractions,
\[
\frac1{v_T^2-u^2}=\frac1{2v_T}\left(\frac1{v_T-u}+\frac1{v_T+u}\right),
\qquad
\int\frac{du}{v_T^2-u^2}=\frac1{2v_T}\ln\left|\frac{v_T+u}{v_T-u}\right|.
\]
Integrating from the nonzero initial speed $v_i$ at $t=0$ therefore gives
\[
\int_{v_i}^{v}\frac{du}{v_T^2-u^2}=\frac bm t,
\qquad
\frac1{2v_T}\ln\left|\frac{v_T+v}{v_T-v}
                         \frac{v_T-v_i}{v_T+v_i}\right|=\frac bm t.
\]
Here $2v_Tb/m=2\sqrt{bg/m}=2/T$, so the logarithm equals $2t/T$.
For $v_i\ne v_T$, put $q=(v_i-v_T)/(v_i+v_T)$. Exponentiating and
inverting the ratio give
\[
\left|\frac{v-v_T}{v+v_T}\right|=|q|\,e^{-2t/T}.
\]
The right side never vanishes, so $v\ne v_T$ at every $t$; by continuity
$v-v_T$ keeps the sign of $v_i-v_T$, which is the sign of $q$. (Since
$v'=(b/m)(v_T^2-v^2)$, the speed moves monotonically from $v_i\ge0$
toward $v_T$, so $v+v_T>0$.) Hence
\[
\frac{v-v_T}{v+v_T}=q e^{-2t/T},\qquad
v(1-qe^{-2t/T})=v_T(1+qe^{-2t/T}).
\]
Divide by $1-qe^{-2t/T}$ and multiply numerator and denominator by
$v_i+v_T$. This gives
\[
\boxed{v(t)=v_T
\frac{(v_i+v_T)+(v_i-v_T)e^{-2t/T}}
     {(v_i+v_T)-(v_i-v_T)e^{-2t/T}}}
=v_T\frac{v_i+v_T\tanh(t/T)}{v_T+v_i\tanh(t/T)}.
\]
For the second form write the numerator as
$v_i(1+e^{-2t/T})+v_T(1-e^{-2t/T})$ and the denominator as
$v_i(1-e^{-2t/T})+v_T(1+e^{-2t/T})$, divide both by $1+e^{-2t/T}$, and
use $(1-e^{-2u})/(1+e^{-2u})=\tanh u$.
For $v_i=v_T$, separation divided by zero; direct substitution instead
gives the equilibrium $v(t)=v_T$, which the boxed formula also includes.
For every $v_i\ge0$ and $t\ge0$, $-1\le q<1$, so its denominator is
positive. It gives $v(0)=v_i$ and approaches $v_T$. For the example's
$m=70\ \mathrm{kg}$, $b=700\ \mathrm{kg/m}$, $g=9.8\ \mathrm{m/s^2}$,
and with $1\ \mathrm{mi/h}=1609.344\ \mathrm m/3600\ \mathrm s
=0.44704\ \mathrm{m/s}$,
\[
\begin{aligned}
v_i&=60(0.44704)=26.8224\ \mathrm{m/s},\\
v_T&=\sqrt{0.98}=0.9899494936\ldots\ \mathrm{m/s},\\
T&=1/\sqrt{98}=0.1010152544\ldots\ \mathrm s.
\end{aligned}
\]
\erratum{Here $v_i>v_T$. The Instructor's Manual's real expression
involving $\operatorname{arctanh}(v_i/v_T)$ is therefore inapplicable,
since $\operatorname{arctanh}w$ is real only for $|w|<1$. The equivalent
real hyperbolic form for this case is}
\[
\erratum{v(t)=v_T\coth\left(t/T+\operatorname{arccoth}(v_i/v_T)\right)},
\quad \operatorname{arccoth}w=\tfrac12\ln\frac{w+1}{w-1}\quad(w>1).
\]
Indeed, with $\phi=\operatorname{arccoth}(v_i/v_T)$ and $u=t/T$, the
addition formula $\coth(u+\phi)=(\coth u\coth\phi+1)/(\coth u+\coth\phi)$,
multiplied above and below by $v_T\tanh u$, gives
$v_T\coth(u+\phi)=v_T(v_i+v_T\tanh u)/(v_T+v_i\tanh u)$, the
$\tanh$ form above.
Differentiating either real form verifies $v'=g(1-v^2/v_T^2)$, which is
$mv'=mg-bv^2$ divided by $m$; e.g., $\operatorname{csch}^2=\coth^2-1$ and
$v_T/T=g$ give
$v'=-(v_T/T)\operatorname{csch}^2(u+\phi)=g(1-v^2/v_T^2)$. The
speed initially decreases, as it must when the parachute opens above
terminal speed. Numerically, $v_i/v_T=27.0947\ldots$, so
$\phi=\frac12\ln[(v_i/v_T+1)/(v_i/v_T-1)]=0.0369243\ldots$, and
$1/T=\sqrt{98}\ \mathrm{s^{-1}}=9.89949\ldots\ \mathrm{s^{-1}}$; with $t$
in seconds,
\[
v(t)\approx0.989949\coth(9.89949\,t+0.0369243)\ \mathrm{m/s}.
\]
For example, $v(0.1\ \mathrm s)=1.2812\ldots\ \mathrm{m/s}$, already
within $0.3\ \mathrm{m/s}$ of $v_T$.

(b) With the free-fall drag coefficient $b=0.25\ \mathrm{kg/m}$,
\[
\boxed{v_T=\sqrt{\frac{70(9.8)}{0.25}}=\sqrt{2744}
=52.38320341\ldots\ \mathrm{m/s}
=188.5795\ldots\ \mathrm{km/h}},
\]
where $1\ \mathrm{m/s}=3.6\ \mathrm{km/h}$.

\Needspace{5\baselineskip}
\par\bigskip
\noindent\textbf{7.2.17}\quad Solve the ODE
\[
(xy^{2}-y)\,dx+x\,dy=0.
\]
\solution
On $x\ne0$, divide by $x$:
\[
y'-\frac yx=-y^2.
\]
The solution $y=0$ is immediate. For $y\ne0$, this Bernoulli equation
becomes linear for $u=1/y$: then $u'=-y'/y^2$, and multiplication by
$-y^{-2}$ gives $-y'/y^2+1/(xy)=1$, i.e.,
\[
u'+\frac ux=1,\qquad (xu)'=xu'+u=x.
\]
Integrating yields $xu=x^2/2+C$, hence $u=(x^2+2C)/(2x)$ and $y=1/u$:
\[
\boxed{y(x)=\frac{2x}{x^2+K}},\qquad \boxed{y\equiv0},
\]
where $K=2C$. Substitution into the original, undivided equation, which
is $xy'=y-xy^2$, checks it: by the quotient rule
$y'=[2(x^2+K)-4x^2]/(x^2+K)^2=2(K-x^2)/(x^2+K)^2$, and
\[
x y'=\frac{2x(K-x^2)}{(x^2+K)^2}
=\frac{2x(x^2+K)-4x^3}{(x^2+K)^2}
=y-xy^2.
\]
The rational solution is valid wherever its denominator is nonzero,
including its smooth extension through $x=0$ when $K\ne0$. The equation
itself forces $y(0)=0$ if a solution is defined there. For $K=0$,
$y=2/x$ has a pole at zero and is used on either side separately.

\Needspace{5\baselineskip}
\par\bigskip
\noindent\textbf{7.2.18}\quad Solve the ODE
\[
(x^{2}-y^{2}e^{y/x})\,dx+(x^{2}+xy)e^{y/x}\,dy=0.
\]
\emph{Hint.}~Note that the quantity $y/x$ in the exponents is of combined
degree zero and does not affect the determination of homogeneity.
\solution
The original equation is defined on $x\ne0$. Set $u=y/x$, so
$y=xu$ and $dy=u\,dx+x\,du$. Substitution gives
\begin{align*}
0&=x^2(1-u^2e^u)dx+x^2(1+u)e^u(u\,dx+x\,du)\\
 &=x^2(1+ue^u)dx+x^3(1+u)e^u du,
\end{align*}
because the coefficient of $dx$ is
$x^2[1-u^2e^u+(u+u^2)e^u]=x^2(1+ue^u)$.
After division by $x^2$ this is already exact, since
$d(ue^u)=(1+u)e^udu$ and the product rule give
\[
(1+ue^u)dx+x(1+u)e^u du=d[x(1+ue^u)]=0.
\]
Therefore $x(1+ue^u)=C$; with $xu=y$ this reads
\[
\boxed{x+ye^{y/x}=C}.
\]
A direct check avoids division by $1+ue^u$ or $1+u$. Using
$d(y/x)=(x\,dy-y\,dx)/x^2$,
\begin{align*}
d(x+ye^{y/x})&=dx+e^u dy+ye^u\frac{x\,dy-y\,dx}{x^2}\\
&=(1-u^2e^u)dx+(1+u)e^u dy,
\end{align*}
which is exactly the original differential form divided by $x^2$.
For real $u$, $ue^u\ge-1/e$ (its minimum, at $u=-1$, where
$(ue^u)'=(1+u)e^u$ vanishes), so $1+ue^u>0$. A constant-ratio solution
$u=\text{const}$ would need $(1+ue^u)dx=0$ in the reduced form, so none
was lost. At $u=-1$ the coefficient of $dy$ vanishes, and the level curve
can have a vertical tangent; the implicit differential-form solution
remains valid there.

\Needspace{8\baselineskip}
\subsection*{7.3 ODEs with Constant Coefficients}

Find the general solutions to the following ODEs. Write the solutions in forms
that are entirely real (i.e., that contain no complex quantities).

\Needspace{5\baselineskip}
\par\bigskip
\noindent\textbf{7.3.1}\quad $y'''-2y''-y'+2y=0$.
\solution
Substitute $y=e^{rx}$ and cancel the nonzero exponential. The
characteristic polynomial factors by grouping:
\[
r^3-2r^2-r+2=r^2(r-2)-(r-2)=(r-2)(r^2-1)=(r-2)(r-1)(r+1).
\]
Its three distinct real roots give
\[
\boxed{y=C_1e^{2x}+C_2e^x+C_3e^{-x}}.
\]
Each exponential solves the ODE because its exponent is a root.
In the Wronskian, the $k$th derivative of $e^{rx}$ is $r^ke^{rx}$, so
the factor $e^{rx}$ comes out of each column:
\[
W=e^{(2+1-1)x}\det\begin{pmatrix}1&1&1\\2&1&-1\\4&1&1\end{pmatrix}
=e^{2x}[(1+1)-(2+4)+(2-4)]=-6e^{2x}\ne0,
\]
expanding along the first row. So these three solutions are
independent and span the three-dimensional solution space.

\Needspace{5\baselineskip}
\par\bigskip
\noindent\textbf{7.3.2}\quad $y'''-2y''+y'-2y=0$.
\solution
For $y=e^{rx}$ the characteristic equation is
\[
r^3-2r^2+r-2=r^2(r-2)+(r-2)=(r-2)(r^2+1)=0.
\]
The roots are $2,i,-i$. Euler's identity $e^{\pm ix}=\cos x\pm i\sin x$
converts the last pair into the real solutions
$(e^{ix}+e^{-ix})/2=\cos x$ and $(e^{ix}-e^{-ix})/(2i)=\sin x$, giving
\[
\boxed{y=C_1e^{2x}+C_2\cos x+C_3\sin x}.
\]
Direct differentiation verifies all three terms. Their Wronskian at
$x=0$ (rows: the values, first and second derivatives of
$e^{2x},\cos x,\sin x$) is, expanding along the first row,
\[
\det\begin{pmatrix}1&1&0\\2&0&1\\4&-1&0\end{pmatrix}
=1\cdot(0+1)-1\cdot(0-4)=5\ne0.
\]
By Abel's formula a Wronskian of solutions is either identically zero or
never zero, so the three solutions are independent and the three arbitrary
real constants provide the general real solution.

\Needspace{5\baselineskip}
\par\bigskip
\noindent\textbf{7.3.3}\quad $y'''-3y'+2y=0$.
\solution
The exponential trial gives the cubic $r^3-3r+2$, which vanishes at
$r=1$; dividing by $r-1$ and factoring the quadratic,
\[
r^3-3r+2=(r-1)(r^2+r-2)=(r-1)^2(r+2).
\]
The repeated root $r=1$ needs a second solution in addition to $e^x$.
To see its form without assuming a multiplicity rule, set $y=e^xu$ and
write $D=d/dx$. Then $(D-1)(e^xu)=e^xu'$ and
$(D+2)(e^xu)=e^x(u'+3u)$; applying $D-1$ twice to $e^x(u'+3u)$
differentiates the factor $u'+3u$ twice, so
\[
(D-1)^2(D+2)y=e^x(u^{(3)}+3u'')=0.
\]
Taking $u=1$ and $u=x$ supplies $e^x$ and $xe^x$; the remaining solution
$u=e^{-3x}$ gives $e^{-2x}$, from the root $r=-2$. Hence
\[
\boxed{y=(C_1+C_2x)e^x+C_3e^{-2x}}.
\]
At $x=0$ the values, first and second derivatives of
$e^x,xe^x,e^{-2x}$ are $(1,0,1)$, $(1,1,-2)$, $(1,2,4)$, so the
Wronskian at $0$ is
$\det\bigl(\begin{smallmatrix}1&0&1\\1&1&-2\\1&2&4\end{smallmatrix}\bigr)
=1\cdot(4+4)+1\cdot(2-1)=9$.
Since the ODE has no $y''$ term, Abel's formula makes the Wronskian
constant, so it equals $9\ne0$ everywhere, which checks independence and
completeness.

\Needspace{5\baselineskip}
\par\bigskip
\noindent\textbf{7.3.4}\quad $y''+2y'+2y=0$.
\solution
The characteristic equation is
$r^2+2r+2=(r+1)^2+1=0$, with roots $-1\pm i$.
Equivalently, substituting $y=e^{-x}u$ gives
\[
y'=e^{-x}(u'-u),\quad y''=e^{-x}(u''-2u'+u),\quad
 y''+2y'+2y=e^{-x}(u''+u).
\]
Thus $u''+u=0$, so $u=C_1\cos x+C_2\sin x$, and
\[
\boxed{y=e^{-x}(C_1\cos x+C_2\sin x)}.
\]
The same real functions are $e^{-x}$ times the real and imaginary parts
of $e^{ix}$ in $e^{(-1\pm i)x}=e^{-x}(\cos x\pm i\sin x)$.
For a common factor $f$, $W(fu_1,fu_2)=f^2W(u_1,u_2)$ (the terms with
$f'$ cancel), so the Wronskian of the two displayed real functions is
$e^{-2x}(\cos^2x+\sin^2x)=e^{-2x}$, which is nonzero for every real $x$.

\Needspace{8\baselineskip}
\subsection*{7.4 Second-Order Linear ODEs}

\Needspace{5\baselineskip}
\par\bigskip
\noindent\textbf{7.4.1}\quad Show that Legendre's equation has regular
singularities at $x=-1$ and $1$, and \erratum{at $\infty$ when
$l(l+1)\ne0$. Classify $\infty$ also when $l=0$ or $l=-1$.}
\solution
Write $\lambda=l(l+1)$ and divide Legendre's equation by $1-x^2$:
\[
y''+P(x)y'+Q(x)y=0,\qquad
P=-\frac{2x}{1-x^2},\quad Q=\frac\lambda{1-x^2}.
\]
At $x=1$, $(x-1)P=2x/(x+1)$ and
$(x-1)^2Q=-\lambda(x-1)/(x+1)$ are analytic. At $x=-1$,
$(x+1)P=2x/(x-1)$ and
$(x+1)^2Q=-\lambda(x+1)/(x-1)$ are analytic.
The simple poles of $P$ are nonremovable, so both are regular singular
points for every $\lambda$.

At infinity put $z=1/x$ and $Y(z)=y(1/z)$. Since $dz/dx=-1/x^2=-z^2$,
the chain rule gives
\[
y'=-z^2Y',\qquad y''=-z^2\frac d{dz}(-z^2Y')=z^4Y''+2z^3Y'.
\]
Substitution into the original equation, with $1-x^2=(z^2-1)/z^2$ and
$-2xy'=2zY'$, gives
\[
(z^2-1)(z^2Y''+2zY')+2zY'+\lambda Y
=z^2(z^2-1)Y''+2z^3Y'+\lambda Y=0,
\]
and division by $z^2(z^2-1)$ gives
\[
Y''+\frac{2z}{z^2-1}Y'+\frac\lambda{z^2(z^2-1)}Y=0.
\]
Thus $z\widetilde P=2z^2/(z^2-1)$ and
$z^2\widetilde Q=\lambda/(z^2-1)$ are analytic at zero. For
$\lambda\ne0$ the double pole of $\widetilde Q$ makes infinity regular
singular.
\erratum{The original unqualified question has an exception: if
$l=0$ or $l=-1$, then $\lambda=0$ and both transformed coefficients
are analytic at $z=0$; infinity is then an ordinary point.}

\Needspace{5\baselineskip}
\par\bigskip
\noindent\textbf{7.4.2}\quad Show that Laguerre's equation, like the Bessel
equation, has a regular singularity at $x=0$ and an irregular singularity at
$x=\infty$.
\solution
For Laguerre's equation $xy''+(1-x)y'+ny=0$,
\[
P=\frac1x-1,\qquad Q=\frac nx.
\]
At $x=0$, $xP=1-x$ and $x^2Q=nx$ are analytic, whereas $P$ has a
simple pole; zero is regular singular.
To analyze infinity, put $z=1/x$, so $y'=-z^2Y'$ and
$y''=z^4Y''+2z^3Y'$ (as in 7.4.1). A normalized ODE
$y''+P(x)y'+Q(x)y=0$ becomes
$z^4Y''+2z^3Y'-z^2P(1/z)Y'+Q(1/z)Y=0$, and dividing by $z^4$ gives
\[
Y''+\left[\frac2z-\frac{P(1/z)}{z^2}\right]Y'
       +\frac{Q(1/z)}{z^4}Y=0.
\]
Here $P(1/z)=z-1$ and $Q(1/z)=nz$, so the transformed coefficients are
\[
\widetilde P=\frac2z-\frac{z-1}{z^2}=\frac1z+\frac1{z^2},\qquad
\widetilde Q=\frac{nz}{z^4}=\frac n{z^3}.
\]
Since $z\widetilde P=1+1/z$ is singular, infinity is irregular even
when $n=0$.

For comparison, Bessel's equation has $P=1/x$, $Q=1-\nu^2/x^2$.
At zero, $xP=1$ and $x^2Q=x^2-\nu^2$ are analytic. At infinity,
$P(1/z)=z$ and $Q(1/z)=1-\nu^2z^2$, so
\[
\widetilde P=\frac2z-\frac z{z^2}=\frac1z,\qquad
\widetilde Q=\frac{1-\nu^2z^2}{z^4}=\frac1{z^4}-\frac{\nu^2}{z^2},
\]
and $z^2\widetilde Q=1/z^2-\nu^2$ is singular. These establish the
same classification for Bessel's equation.

\Needspace{5\baselineskip}
\par\bigskip
\noindent\textbf{7.4.3}\quad Show that Chebyshev's equation, like the Legendre
equation \erratum{when $l(l+1)\ne0$}, has regular singularities at $x=-1$, $1$, and $\infty$.
\solution
Chebyshev's equation is $(1-x^2)y''-xy'+n^2y=0$, so
\[
P=-\frac{x}{1-x^2},\qquad Q=\frac{n^2}{1-x^2}.
\]
At $1$, $(x-1)P=x/(x+1)$ and
$(x-1)^2Q=-n^2(x-1)/(x+1)$ are analytic; at $-1$ the corresponding
expressions are $x/(x-1)$ and $-n^2(x+1)/(x-1)$.
The poles of $P$ ensure both points are regular singular.

With $z=1/x$,
$P(1/z)=-(1/z)/(1-z^{-2})=z/(1-z^2)$ and
$Q(1/z)=n^2/(1-z^{-2})=n^2z^2/(z^2-1)$, so the transformed coefficients
of 7.4.2 are
\[
\widetilde P=\frac2z-\frac1{z(1-z^2)}
=\frac{1-2z^2}{z(1-z^2)},\qquad
\widetilde Q=\frac{n^2}{z^2(z^2-1)}.
\]
Then $z\widetilde P=(1-2z^2)/(1-z^2)$ and
$z^2\widetilde Q=n^2/(z^2-1)$ are analytic at zero.
Since $z\widetilde P\to1$ as $z\to0$, $\widetilde P$ always has a simple
pole, and infinity is regular singular even for $n=0$.
\erratum{The Instructor's Manual's first displayed normalized Chebyshev
equation has
$+x/(1-x^2)$ multiplying $y'$. Division of the original equation gives
$-x/(1-x^2)$ instead; its transformed $P$ coefficient must therefore
also be the one calculated above. The singularity classification is
unchanged.}

\Needspace{5\baselineskip}
\par\bigskip
\noindent\textbf{7.4.4}\quad Show that Hermite's equation has no singularity
other than an irregular singularity at $x=\infty$.
\solution
In $y''-2xy'+2\alpha y=0$, the normalized coefficients
$P=-2x$ and $Q=2\alpha$ are polynomials. They have no singularities at
any finite complex $x$.
For $z=1/x$, $P(1/z)=-2/z$ and $Q(1/z)=2\alpha$, so the transformation
of 7.4.2 gives $\widetilde P=2/z+(2/z)/z^2$, $\widetilde Q=2\alpha/z^4$,
i.e., the transformed equation is
\[
Y''+\left(\frac2z+\frac2{z^3}\right)Y'
       +\frac{2\alpha}{z^4}Y=0.
\]
Multiplying the coefficient of $Y'$ by $z$ gives $2+2/z^2$, which
is not analytic or finite at zero. Thus infinity is an irregular singular
point for every $\alpha$, including $\alpha=0$.

\Needspace{5\baselineskip}
\par\bigskip
\noindent\textbf{7.4.5}\quad Show that the substitution
\[
x\to\frac{1-x}{2},\quad a=-l,\quad b=l+1,\quad c=1
\]
converts the hypergeometric equation into Legendre's equation.
\erratum{Use $t(t-1)u''+[(1+a+b)t-c]u'+ab\,u=0$;
Table 7.1 has the wrong sign on $c$.}
\solution
Distinguish the hypergeometric variable $t$ from the new variable $x$:
\[
t(t-1)u_{tt}+[(1+a+b)t-c]u_t+ab\,u=0,
\qquad t=\frac{1-x}{2},\quad y(x)=u(t).
\]
\erratum{The coefficient is $(1+a+b)t-c$, as in Exercise 7.5.5.
Table 7.1 of the textbook prints $+c$ in this sign convention; with it
the coefficient of $u_t$ would be $2t+1=2-x$, and the steps below would
give $(1-x^2)y''+(4-2x)y'+l(l+1)y=0$, which is not Legendre's equation.}
Since $x=1-2t$, $dx/dt=-2$, and the chain rule gives $u_t=-2y'$ and
$u_{tt}=4y''$. Also $t-1=-(1+x)/2$, so
$t(t-1)=-(1-x)(1+x)/4=(x^2-1)/4$. With $a=-l$, $b=l+1$, $c=1$,
\[
1+a+b=2,\qquad 2t-1=-x,\qquad ab=-l(l+1).
\]
Substitution therefore gives
\[
(x^2-1)y''+(-x)(-2y')-l(l+1)y=0.
\]
Multiplying by $-1$ yields exactly
\[
\boxed{(1-x^2)y''-2xy'+l(l+1)y=0}.
\]

\Needspace{8\baselineskip}
\subsection*{7.5 Series Solutions---Frobenius' Method}

\Needspace{5\baselineskip}
\par\bigskip
\noindent\textbf{7.5.1}\quad \emph{Uniqueness theorem.} The function $y(x)$
satisfies a second-order, linear, homogeneous differential equation. At
$x=x_{0}$, $y(x)=y_{0}$ and $dy/dx=y'_{0}$. Show that $y(x)$ is unique, in
that no other solution of this differential equation passes through the points
$(x_{0},y_{0})$ with a slope of $y'_{0}$.
\erratum{Assume a connected interval on which the normalized coefficients
are continuous. For the Taylor-series hint, assume they are analytic
and $x_0$ is an ordinary point.}

\emph{Hint.}~Assume a second solution satisfying these conditions and compare
the Taylor series expansions.
\solution
\erratum{Uniqueness is guaranteed when the normalized coefficients are
continuous near the initial point (in particular at an ordinary point);
at a singular point it can fail. The hint using Taylor series additionally
requires analytic coefficients. The unqualified claim, and the
Instructor's Manual's extension to regular singular points, are false.}
For example, $xy''-y'=0$ (regular singular at $0$) has the family
$y=Cx^2$, since $x(2C)-2Cx=0$; all members have $y(0)=y'(0)=0$, although
they differ elsewhere.

For the Taylor proof, write $y''+P(x)y'+Q(x)y=0$, with $P,Q$ analytic
at $x_0$. If $y$ and $\widetilde y$ share the initial data, their
difference $u=y-\widetilde y$ solves the same equation (it is linear and
homogeneous) and has $u(x_0)=u'(x_0)=0$. The equation gives
$u''(x_0)=-P(x_0)u'(x_0)-Q(x_0)u(x_0)=0$.
Differentiating $k$ times and applying the product rule gives
\[
u^{(k+2)}=-\sum_{j=0}^k\binom kj
 \{P^{(j)}u^{(k-j+1)}+Q^{(j)}u^{(k-j)}\}.
\]
The right side involves derivatives of $u$ only up to order $k+1$ (the
term $j=0$). If all derivatives through order $k+1$ vanish at $x_0$, the
right side vanishes there, and so does $u^{(k+2)}(x_0)$. Induction on $k$
therefore makes every Taylor coefficient of $u$ zero. The hint
presupposes that $u$, like the series solutions constructed in
Section~7.5, is analytic at the ordinary point $x_0$, i.e., equal to its
convergent Taylor series there. This is true, but Section~7.5 only
constructs series solutions; that \emph{every} solution is analytic is
normally derived from the uniqueness theorem itself, so the integral
estimate below is the self-contained proof. Granting the hint's
assumption, $u=0$ in a neighborhood of $x_0$. The same argument
applies at every point of the connected ordinary interval: the set where $u=u'=0$ is closed (by continuity) and open (by
this argument), and it contains $x_0$, so it is the whole interval. This
proves uniqueness throughout it.

For merely continuous $P,Q$, one can replace the Taylor argument by an
integral estimate. With zero initial data, integrating $u'$ and
$u''=-Pu'-Qu$ from $x_0$ gives
\[
u(x)=\int_{x_0}^xu'(t)dt,\qquad
u'(x)=-\int_{x_0}^x[P(t)u'(t)+Q(t)u(t)]dt.
\]
On a compact interval $[x_0,x_1]$ to the right, $P$ and $Q$ are
continuous, hence bounded, so we can choose a constant $K$ with
$K\ge1+|P(t)|+|Q(t)|$ for all $t\in[x_0,x_1]$, and put
$F(x)=\int_{x_0}^x(|u|+|u'|)dt$. The two integrals give
$|u(x)|\le F(x)$ and
$|u'(x)|\le\int_{x_0}^x(|P||u'|+|Q||u|)dt\le(K-1)F(x)$, so
$F'=|u|+|u'|\le KF$, and $F(x_0)=0$.
Thus $(e^{-K(x-x_0)}F)'=e^{-K(x-x_0)}(F'-KF)\le0$, so
$e^{-K(x-x_0)}F(x)\le F(x_0)=0$, while $F\ge0$; this forces $F=0$, hence
$F'=|u|+|u'|=0$ and
$u=u'=0$. Reverse the integration direction for points to the left.
This establishes the continuous-coefficient version without assuming that
a smooth function equals its Taylor series.

\Needspace{5\baselineskip}
\par\bigskip
\noindent\textbf{7.5.2}\quad A series solution of Eq.~(7.23) is attempted,
expanding about the point $x=x_{0}$. If $x_{0}$ is an ordinary point, show
that the indicial equation has roots $s=0,1$.
\solution
Set $t=x-x_0$ and write the normalized equation as
$y''+P(t)y'+Q(t)y=0$, where $P,Q$ are analytic at $t=0$.
For a Frobenius trial
\[
y=\sum_{j=0}^\infty a_jt^{s+j},\quad a_0\ne0,
\quad y'=\sum_j(s+j)a_jt^{s+j-1},\quad
 y''=\sum_j(s+j)(s+j-1)a_jt^{s+j-2},
\]
the lowest possible power is $t^{s-2}$. Only $y''$ contributes to it:
since $P(t)=P(0)+O(t)$ and $Q(t)=Q(0)+O(t)$ have no negative powers,
$P y'$ starts at $t^{s-1}$ and $Qy$ at $t^s$. Hence
\[
a_0s(s-1)=0,\qquad \boxed{s=0\ \hbox{or}\ s=1}.
\]
These represent a solution that can be nonzero at $x_0$ and one with a
simple zero there. The coefficient of $t^{s-1}$ receives $(s+1)sa_1$
from $y''$ (term $j=1$) and $P(0)sa_0$ from $Py'$, so it is
$s(s+1)a_1+sP(0)a_0=0$; for $s=1$ this gives
$a_1=-P(0)a_0/2$, so ordinary-point exponents alone do not generally
force $a_1$ to vanish.

\Needspace{5\baselineskip}
\par\bigskip
\noindent\textbf{7.5.3}\quad In the development of a series solution of the
simple harmonic oscillator (SHO) equation, the second series coefficient
$a_{1}$ was neglected except to set it equal to zero. From the coefficient of
the next-to-the-lowest power of $x$, $x^{s-1}$, develop a second-indicial
type equation.
\begin{enumerate}
\item[(a)] (SHO equation with $s=0$). Show that $a_{1}$ may be assigned any
  finite value (including zero).
\item[(b)] (SHO equation with $s=1$). Show that $a_{1}$ must be set equal to
  zero.
\end{enumerate}
\solution
Use $y''+\omega^2y=0$ and
$y=\sum_{j\ge0}a_jx^{s+j}$ with $a_0\ne0$. Substitution gives
\[
\sum_{j\ge0}(s+j)(s+j-1)a_jx^{s+j-2}
 +\omega^2\sum_{j\ge0}a_jx^{s+j}=0.
\]
The coefficients of $x^{s-2}$ and $x^{s-1}$ come from the first sum only
($j=0$ and $j=1$); they are respectively
\[
s(s-1)a_0=0,\qquad \boxed{s(s+1)a_1=0}.
\]
Replacing $j$ by $j+2$ in the first sum, the coefficient of $x^{s+j}$,
$j\ge0$, gives
\[
(s+j+2)(s+j+1)a_{j+2}+\omega^2a_j=0.
\]
(a) For $s=0$ the second equation is $0=0$, leaving $a_1$ arbitrary.
The subsequent equations give
\[
a_{j+2}=-\frac{\omega^2a_j}{(j+2)(j+1)},
\]
so the even coefficients depend on $a_0$ and the odd coefficients on
$a_1$: iterating, $a_{2k}=(-\omega^2)^ka_0/(2k)!$ and
$a_{2k+1}=(-\omega^2)^ka_1/(2k+1)!$. In particular, comparison with the
Maclaurin series of $\cos\omega x$ and $\sin\omega x$ gives
$y=a_0\cos\omega x+(a_1/\omega)\sin\omega x$ for $\omega\ne0$.

(b) For $s=1$ the second equation is $2a_1=0$, hence $a_1=0$.
The recurrence becomes $a_{j+2}=-\omega^2a_j/[(j+3)(j+2)]$, so
$a_{2k}=(-\omega^2)^ka_0/(2k+1)!$ and the series
$a_0\sum_k(-1)^k\omega^{2k}x^{2k+1}/(2k+1)!$ is
$a_0\sin(\omega x)/\omega$. Thus setting $a_1=0$ in the
$s=0$ series is harmless only when the separate $s=1$ solution is also
retained. For $\omega=0$ the two solutions reduce to $1,x$.

\Needspace{5\baselineskip}
\par\bigskip
\noindent\textbf{7.5.4}\quad Analyze the series solutions of the following
differential equations to see when $a_{1}$ may be set equal to zero without
irrevocably losing anything and when $a_{1}$ must be set equal to zero.
\begin{enumerate}
\item[(a)] Legendre,\quad(b) Chebyshev,\quad(c) Bessel,\quad(d) Hermite.
\end{enumerate}
\par\nobreak\hfill\emph{ANS.}\enspace
(a) Legendre, (b) Chebyshev, and (d) Hermite: For $s=0$, $a_{1}$ may be set
equal to zero; for $s=1$, $a_{1}$ must be set equal to zero.\par
\nobreak\hfill(c) Bessel: $a_{1}$ must be set equal to zero (except for
$s=\pm n=-\frac{1}{2}$).
\solution
Insert $y=\sum_{j\ge0}a_jx^{s+j}$, $a_0\ne0$. For Legendre,
Chebyshev, and Hermite the coefficient multiplying $y'$ starts with $x$,
so their lowest two powers come only from $y''$:
\[
s(s-1)a_0=0,\qquad s(s+1)a_1=0.
\]
The coefficient of $x^{s+j}$ receives $(s+j+2)(s+j+1)a_{j+2}$ from
$y''$ and multiples of $a_j$ from the other terms:
$-(s+j)(s+j-1)-2(s+j)+l(l+1)=-(s+j)(s+j+1)+l(l+1)$ from
$-x^2y''-2xy'+l(l+1)y$ (Legendre),
$-(s+j)(s+j-1)-(s+j)+n^2=-(s+j)^2+n^2$ from $-x^2y''-xy'+n^2y$
(Chebyshev), and $-2(s+j)+2\alpha$ from $-2xy'+2\alpha y$ (Hermite).
Setting the coefficient to zero gives the explicit recurrences,
respectively,
\begin{align*}
\text{Legendre: }\quad
 a_{j+2}&=\frac{(s+j)(s+j+1)-l(l+1)}{(s+j+2)(s+j+1)}a_j,\\
\text{Chebyshev: }\quad
 a_{j+2}&=\frac{(s+j)^2-n^2}{(s+j+2)(s+j+1)}a_j,\\
\text{Hermite: }\quad
 a_{j+2}&=\frac{2(s+j-\alpha)}{(s+j+2)(s+j+1)}a_j.
\end{align*}
Thus (a), (b), and (d) have arbitrary $a_1$ for $s=0$, but require
$a_1=0$ for $s=1$. Each recurrence links $a_{j+2}$ to $a_j$ only, so it
preserves parity. Moreover, each depends on $s$ and $j$ only through the
power $s+j$; hence the odd coefficients $a_1,a_3,\ldots$ of the $s=0$
branch obey the same recurrence as $a_0,a_2,\ldots$ of the $s=1$ branch,
and the $s=0$ series with $a_1\ne0$ is $a_0$ times the even solution plus
$a_1$ times the $s=1$ solution. Choosing $a_1=0$
in the $s=0$ branch gives the even solution; the omitted odd branch is
then supplied separately by $s=1$. It would be lost if only the even
branch were retained.

(c) Bessel's equation is $x^2y''+xy'+(x^2-n^2)y=0$. At $x^{s+j}$ the
first two terms contribute $[(s+j)(s+j-1)+(s+j)]a_j=(s+j)^2a_j$ and
$x^2y$ contributes $a_{j-2}$. Its lowest two
coefficients and subsequent recurrence are therefore
\[
(s^2-n^2)a_0=0,\quad [(s+1)^2-n^2]a_1=0,\quad
[(s+j)^2-n^2]a_j+a_{j-2}=0\quad(j\ge2).
\]
At either indicial root $s^2=n^2$, the second coefficient reduces to
$(s+1)^2-n^2=s^2+2s+1-n^2=2s+1$, so $(2s+1)a_1=0$. Hence $a_1$ is
forced to zero except when
\[
\boxed{s=-\tfrac12,\qquad n^2=\tfrac14}.
\]
In that exceptional lower-root branch, $a_1$ is free and generates the
other, $s=+1/2$, solution. This conclusion concerns $a_1$ only; for
integer orders a later vanishing denominator in the lower-root recurrence
can require a logarithmic second solution.

\Needspace{5\baselineskip}
\par\bigskip
\noindent\textbf{7.5.5}\quad Obtain a series solution of the hypergeometric
equation
\[
x(x-1)y''+[(1+a+b)x-c]y'+aby=0.
\]
Test your solution for convergence.
\solution
Multiply the equation by $-1$ to use
$x(1-x)y''+[c-(a+b+1)x]y'-aby=0$.
Insert $y=\sum_{j\ge0}A_jx^{s+j}$, $A_0\ne0$. The contributions from
$x y''+c y'$ at power $x^{s+j-1}$ are
$[(s+j)(s+j-1)+c(s+j)]A_j=(s+j)(s+j+c-1)A_j$; those from
$-x^2y''-(a+b+1)xy'-aby$ at power $x^{s+j}$ are
\[
-[(s+j)(s+j-1)+(a+b+1)(s+j)+ab]A_j=-(s+j+a)(s+j+b)A_j,
\]
because the bracket equals $(s+j)^2+(a+b)(s+j)+ab$. The lowest power
$x^{s-1}$ gives the indicial equation; replacing $j$ by $j+1$ in the
first group, the coefficient of $x^{s+j}$ gives
$(s+j+1)(s+j+c)A_{j+1}-(s+j+a)(s+j+b)A_j=0$. Thus
\[
s(s+c-1)=0,
\quad
\boxed{A_{j+1}=\frac{(s+j+a)(s+j+b)}{(s+j+1)(s+j+c)}A_j}.
\]
With $(q)_j=q(q+1)\cdots(q+j-1)$ and $(q)_0=1$, iterating from $A_0=1$
for $s=0$ gives
$A_j=\prod_{i=0}^{j-1}\frac{(i+a)(i+b)}{(i+1)(i+c)}=\frac{(a)_j(b)_j}{(c)_jj!}$;
for $s=1-c$ the factors become $(i+a-c+1)(i+b-c+1)/[(i+1)(i+2-c)]$.
The two generic Frobenius solutions are therefore
\begin{align*}
y_1&=\sum_{j=0}^\infty\frac{(a)_j(b)_j}{(c)_j j!}x^j
=1+\frac{ab}{c}x+\frac{a(a+1)b(b+1)}{2c(c+1)}x^2+\cdots,\\
y_2&=x^{1-c}\sum_{j=0}^\infty
\frac{(a-c+1)_j(b-c+1)_j}{(2-c)_j j!}x^j.
\end{align*}
For noninteger $c$ no denominator $(j+1)(j+c)$ or $(j+1)(j+2-c)$
vanishes, and the two solutions are independent: their leading powers are
$1$ and $x^{1-c}$, and with $y_1=1+O(x)$, $y_2=x^{1-c}[1+O(x)]$ their
Wronskian $y_1y_2'-y_1'y_2$ starts as $(1-c)x^{-c}\ne0$.
A branch of the fractional power is fixed when complex $x$ is used.

\erratum{The displayed coefficient formulas require nonzero denominators.
At integer $c$, one cannot automatically assert two independent solutions
of these exact forms.} A valid first series always comes from the root
$s=0$ if $c\ge1$ and $s=1-c$ if $c\le0$: its recurrence denominators
do not vanish. When $c$ is an integer, the roots $0$ and $1-c$ coincide
($c=1$) or differ by a nonzero integer; then the lower root either
repeats the first series or its recurrence meets a zero denominator (at
$j=c-2$ if $c\ge2$, at $j=-c$ if $c\le0$), so the second solution may
require a logarithm. To construct it,
normalize the equation and split the first-derivative coefficient:
\[
P(x)=\frac{c-(a+b+1)x}{x(1-x)}
=\frac cx+\frac{c-a-b-1}{1-x}.
\]
Consequently $e^{-\int Pdx}=x^{-c}(1-x)^{c-a-b-1}$, and reduction of
order (the second-solution formula, Eq.~(7.67) of Section~7.6) gives
\[
y_2=y_1\int^x\frac{t^{-c}(1-t)^{c-a-b-1}}{y_1(t)^2}\,dt.
\]
Expand the analytic factors and integrate each power near zero; a
$t^{-1}$ term produces $\ln t$. On a branch avoiding $0,1$ and zeros
of $y_1$, its Wronskian is
$W(y_1,y_2)=y_1^2\,(y_2/y_1)'=x^{-c}(1-x)^{c-a-b-1}\ne0$.
This supplies the missing independent solution in every integer-$c$
case. If a numerator factor $s+j+a$ or $s+j+b$ vanishes at some $j$
before a denominator factor does, then $A_{j+1}=0$ and the series
terminates; the resulting polynomial (times $x^s$) is a solution,
because the undivided recurrence
$(s+j+1)(s+j+c)A_{j+1}=(s+j+a)(s+j+b)A_j$ then holds for every $j$,
reading $0=0$ beyond the termination, including at a zero denominator.

For a nonterminating series, its successive terms satisfy
\[
\left|\frac{A_{j+1}x^{j+1}}{A_jx^j}\right|
=|x|\left|\frac{(j+s+a)(j+s+b)}{(j+s+1)(j+s+c)}\right|
\longrightarrow |x|.
\]
By the ratio test it converges absolutely for $|x|<1$ and diverges for
$|x|>1$. For boundary convergence put $\sigma=c-a-b$. Writing
$(q)_j=\Gamma(j+q)/\Gamma(q)$, the coefficients of $y_1$ are
\[
A_j=\frac{\Gamma(c)}{\Gamma(a)\Gamma(b)}\,
\frac{\Gamma(j+a)}{\Gamma(j+1)}\,\frac{\Gamma(j+b)}{\Gamma(j+c)}.
\]
For large $j$ use Stirling's series (Section~13.4) with
$\Gamma(z+1)=z\Gamma(z)$:
$\ln\Gamma(z)=(z-\frac12)\ln z-z+\frac12\ln2\pi+O(z^{-1})$. For fixed $u$,
$\ln(j+u)=\ln j+u/j+O(j^{-2})$, so
$(j+u-\frac12)\ln(j+u)=(j+u-\frac12)\ln j+u+O(j^{-1})$ and
\[
\ln\Gamma(j+u)=(j+u-\tfrac12)\ln j-j+\tfrac12\ln2\pi+O(j^{-1}).
\]
Subtracting the same expansion with $v$ leaves $(u-v)\ln j+O(j^{-1})$,
and exponentiating gives
$\Gamma(j+u)/\Gamma(j+v)=j^{u-v}[1+O(j^{-1})]$. With $(u,v)=(a,1)$ and
$(b,c)$ this gives, in a
nonterminating generic case ($a,b,c$ not zero or negative integers),
\[
A_j=Kj^{(a-1)+(b-c)}[1+O(j^{-1})]=Kj^{-\sigma-1}[1+O(j^{-1})],
\qquad K=\frac{\Gamma(c)}{\Gamma(a)\Gamma(b)}\ne0.
\]
Therefore at $x=1$
convergence holds precisely when $\operatorname{Re}\sigma>0$.
Indeed, absolute convergence is then comparison with
$j^{-1-\operatorname{Re}\sigma}$. If $-1<\operatorname{Re}\sigma\le0$,
compare the leading powers $Kj^{-\sigma-1}$ with
$\int_1^Nt^{-\sigma-1}dt=(N^{-\sigma}-1)/(-\sigma)$, or $\ln N$ when
$\sigma=0$. Since
$|d(t^{-\sigma-1})/dt|=|\sigma+1|\,t^{-\operatorname{Re}\sigma-2}$, for
$j\le t\le j+1$ we have
$|j^{-\sigma-1}-t^{-\sigma-1}|\le|\sigma+1|\,j^{-\operatorname{Re}\sigma-2}$,
which is summable. Splitting the integral at the integers,
\[
\sum_{j=1}^{N-1}j^{-\sigma-1}-\int_1^Nt^{-\sigma-1}dt
=\sum_{j=1}^{N-1}\int_j^{j+1}\bigl(j^{-\sigma-1}-t^{-\sigma-1}\bigr)dt
\]
therefore tends to a finite limit. The partial sums thus contain a term
proportional to $N^{-\sigma}$, or $\ln N$, with no limit:
$|N^{-\sigma}|=N^{-\operatorname{Re}\sigma}\to\infty$ if
$\operatorname{Re}\sigma<0$, while $N^{-i\tau}=e^{-i\tau\ln N}$ keeps
circling the unit circle if $\sigma=i\tau\ne0$. The
$O(j^{-\operatorname{Re}\sigma-2})$ correction is summable. For
$\operatorname{Re}\sigma\le-1$ even the individual terms fail to tend
to zero. At $x=-1$, put $p=\sigma+1$. When $\operatorname{Re}p>0$,
$j^{-p}\to0$ and
$(j+1)^{-p}-j^{-p}=O(j^{-\operatorname{Re}p-1})$ is absolutely
summable. Summation by parts with the bounded partial sums of
$(-1)^j$ (Dirichlet's test) therefore proves convergence. The
coefficient correction is
also absolutely summable. Thus the condition there is
$\operatorname{Re}\sigma>-1$, with absolute convergence only for
$\operatorname{Re}\sigma>0$. The second series has the same
$\sigma$, since $(2-c)-(a-c+1)-(b-c+1)=c-a-b$.
For any other point on $|x|=1$ with $x\ne1$, the partial sums of $x^j$
are bounded by $2/|1-x|$, so the same summation-by-parts argument as
at $-1$ gives convergence precisely when $\operatorname{Re}\sigma>-1$.
A terminating power-series factor is a polynomial and converges
for every finite $x$; the second solution still retains its possibly
singular or multivalued prefactor $x^{1-c}$.

\Needspace{5\baselineskip}
\par\bigskip
\noindent\textbf{7.5.6}\quad Obtain two series solutions of the confluent
hypergeometric equation
\[
xy''+(c-x)y'-ay=0.
\]
Test your solutions for convergence.
\erratum{For integer $c$, allow a generalized Frobenius solution containing
a logarithm; two independent pure power-series forms need not exist.}
\solution
Put $y=\sum_{j\ge0}A_jx^{s+j}$, $A_0\ne0$. Then
$x y''+c y'$ contributes
$[(s+j)(s+j-1)+c(s+j)]A_jx^{s+j-1}=(s+j)(s+j+c-1)A_jx^{s+j-1}$,
whereas $-xy'-ay$ contributes $-(s+j+a)A_jx^{s+j}$. Matching powers
(the lowest, $x^{s-1}$, then $x^{s+j}$ after replacing $j$ by $j+1$ in
the first group),
\[
s(s+c-1)=0,\qquad
\boxed{A_{j+1}=\frac{s+j+a}{(s+j+1)(s+j+c)}A_j}.
\]
Iterating from $A_0=1$ gives $A_j=(a)_j/[(c)_jj!]$ for $s=0$; for
$s=1-c$ the factors are $(i+a-c+1)/[(i+1)(i+2-c)]$, $i=0,\ldots,j-1$.
For noninteger $c$ no denominator vanishes, and the independent
solutions, each with leading coefficient one, are
\begin{align*}
y_1&=\sum_{j=0}^\infty\frac{(a)_j}{(c)_j j!}x^j
=1+\frac acx+\frac{a(a+1)}{2c(c+1)}x^2+\cdots,\\
y_2&=x^{1-c}\sum_{j=0}^\infty
\frac{(a-c+1)_j}{(2-c)_j j!}x^j.
\end{align*}
The first is conventionally $M(a,c,x)$; the second is
$x^{1-c}M(a-c+1,2-c,x)$. As in 7.5.5, $y_1=1+O(x)$ and
$y_2=x^{1-c}[1+O(x)]$, so their Wronskian has leading term
$(1-c)x^{-c}$, verifying independence.
For either coefficient sequence,
\[
\left|\frac{A_{j+1}x^{j+1}}{A_jx^j}\right|
=|x|\left|\frac{s+j+a}{(s+j+1)(s+j+c)}\right|
\sim\frac{|x|}{j}\longrightarrow0.
\]
Thus both power-series factors converge for every finite $x$.
The prefactor $x^{1-c}$ may still be singular or multivalued at zero;
convergence of its power-series factor does not remove that singularity.
A terminating power-series factor is a polynomial; the second solution
still includes its displayed power prefactor.

\erratum{For integer $c$, two independent pure Frobenius series need not
exist; a logarithmic solution may be necessary. The ordinary formulas
above must not be used through a zero denominator.}
\par\Needspace{4\baselineskip}
A constructive series prescription for all integer $c$ is as follows. Choose
$s=\max(0,1-c)$ and obtain $y_1=x^sA(x)$ with $A(0)=1$ by the displayed
recurrence; its denominators are nonzero. Write
\[
A(x)^2=\sum_{j\ge0}B_jx^j,\quad B_j=\sum_{r=0}^jA_rA_{j-r},
\quad \frac{e^x}{A(x)^2}=\sum_{j\ge0}h_jx^j.
\]
Since $B_0=1$, comparing the coefficients of $x^k$ in
$A(x)^2\sum_jh_jx^j=e^x$, i.e., $\sum_{j=0}^kB_jh_{k-j}=1/k!$, gives
$h_0=1$ and $h_k=1/k!-\sum_{j=1}^k B_jh_{k-j}$.
The normalized equation has $P=c/x-1$, so $e^{-\int P\,dx}=x^{-c}e^x$,
and reduction of order gives
\[
y_2=y_1\int^x\frac{t^{-c}e^t}{y_1(t)^2}\,dt
=y_1\int^x t^{-c-2s}\frac{e^t}{A(t)^2}\,dt
=y_1\int^x\sum_{k\ge0}h_kt^{k-c-2s}\,dt.
\]
Integrating term by term, $t^{k-c-2s}$ gives
$x^{k-c-2s+1}/(k-c-2s+1)$, except for $k=c+2s-1$, where it gives
$\ln x$. This yields the explicit generalized series
\[
y_2=x^sA(x)\left[
\sum_{\substack{k\ge0\\k\ne c+2s-1}}
 \frac{h_kx^{k-c-2s+1}}{k-c-2s+1}
+h_{c+2s-1}\ln x\right].
\]
Here $c+2s-1$ is a nonnegative integer ($c-1\ge0$ if $c\ge1$, and
$1-c\ge1$ if $c\le0$), and the logarithm is omitted
if its coefficient vanishes. This expansion converges locally, near zero
where $A(x)\ne0$; its Wronskian with $y_1$ is
$y_1^2(y_2/y_1)'=e^xx^{-c}\ne0$.
For example, $a=c=1$ gives $s=0$, $A_j=(1)_j/[(1)_jj!]=1/j!$, so
$y_1=e^x$; then $e^x/A^2=e^{-x}$, $h_k=(-1)^k/k!$, $c+2s-1=0$, and
$y_2=e^x[\ln x+\sum_{k\ge1}(-x)^k/(k k!)]$, displaying the necessary
logarithm explicitly. Real solutions on the negative half-line use
$\ln|x|$ and an appropriate real power convention.

\Needspace{5\baselineskip}
\par\bigskip
\noindent\textbf{7.5.7}\quad A quantum mechanical analysis of the Stark effect
(parabolic coordinates) leads to the differential equation
\[
\frac{d}{d\xi}\left(\xi\frac{du}{d\xi}\right)
+\left(\frac{1}{2}E\xi+\alpha-\frac{m^{2}}{4\xi}-\frac{1}{4}F\xi^{2}\right)u=0.
\]
Here $\alpha$ is a constant, $E$ is the total energy, and $F$ is a constant
such that $Fz$ is the potential energy added to the system by the introduction
of an electric field.

Using the larger root of the indicial equation, develop a power-series
solution about $\xi=0$. Evaluate the first three coefficients in terms of
$a_{0}$.
\par\nobreak\hfill\emph{ANS.}\enspace
Indicial equation $s^{2}-\dfrac{m^{2}}{4}=0$,\par
\nobreak\hfill
\erratum{$M=|m|$,}
$u(\xi)=a_{0}\xi^{\erratum{M}/2}\left[1-\dfrac{\alpha}{\erratum{M}+1}\xi
+\left(\dfrac{\alpha^{2}}{2(\erratum{M}+1)(\erratum{M}+2)}
-\dfrac{E}{4(\erratum{M}+2)}\right)\xi^{2}+\cdots\right]$.\par
Note that the perturbation $F$ does not appear until $a_{3}$ is included.
\solution
\erratum{The larger indicial root is $|m|/2$ for a real azimuthal quantum
number. The printed answer implicitly assumes $m\ge0$; use $M=|m|$
for either sign.}
Expand the derivative and multiply by $\xi$:
\[
\xi^2u''+\xi u'-\frac{m^2}{4}u
 +\alpha\xi u+\frac E2\xi^2u-\frac F4\xi^3u=0.
\]
For $u=\sum_{j\ge0}a_j\xi^{s+j}$, set negative-index coefficients to
zero. At $\xi^{s+j}$, $\xi^2u''+\xi u'$ contributes
$[(s+j)(s+j-1)+(s+j)]a_j=(s+j)^2a_j$, while $\alpha\xi u$,
$(E/2)\xi^2u$ and $-(F/4)\xi^3u$ contribute $\alpha a_{j-1}$,
$(E/2)a_{j-2}$ and $-(F/4)a_{j-3}$. Matching $\xi^{s+j}$ gives
\[
\left[(s+j)^2-\frac{m^2}{4}\right]a_j
 +\alpha a_{j-1}+\frac E2a_{j-2}-\frac F4a_{j-3}=0.
\]
At $j=0$ this yields $s^2-m^2/4=0$, so $s=\pm M/2$; choose the larger
root $s=M/2$. For $j\ge1$ the coefficient of $a_j$ is
$(M/2+j)^2-M^2/4=j(j+M)\ne0$, so
\[
a_j=\frac{-\alpha a_{j-1}-(E/2)a_{j-2}+(F/4)a_{j-3}}{j(j+M)}.
\]
In particular, $j=1$ and $j=2$ give, with $-\alpha a_1=\alpha^2a_0/(M+1)$,
\[
a_1=-\frac\alpha{M+1}a_0,
\quad
a_2=\frac{-\alpha a_1-(E/2)a_0}{2(M+2)}
=a_0\left[\frac{\alpha^2}{2(M+1)(M+2)}-\frac E{4(M+2)}\right].
\]
Therefore the requested result is
\[
\boxed{u=a_0\xi^{M/2}\left[1-\frac\alpha{M+1}\xi
 +\left(\frac{\alpha^2}{2(M+1)(M+2)}-\frac E{4(M+2)}\right)\xi^2
 +\cdots\right]}.
\]
The next coefficient ($j=3$, with $-(E/2)a_1=E\alpha a_0/[2(M+1)]$) is
\[
a_3=\frac1{3(M+3)}\left[-\alpha a_2
 +\frac{E\alpha}{2(M+1)}a_0+\frac F4a_0\right],
\]
so $F$ first enters at $a_3$, as claimed.

\Needspace{5\baselineskip}
\par\bigskip
\noindent\textbf{7.5.8}\quad For the special case of no azimuthal dependence,
the quantum mechanical analysis of the hydrogen molecular ion leads to the
equation
\[
\frac{d}{d\eta}\left[(1-\eta^{2})\frac{du}{d\eta}\right]+\alpha u
+\beta\eta^{2}u=0.
\]
Develop a power-series solution for $u(\eta)$. Evaluate the first three
\erratum{coefficients not forced to vanish by parity ($a_0,a_2,a_4$) in
terms of $a_0$; special parameter values may make some of them zero.}
\par\nobreak\hfill\emph{ANS.}\enspace
Indicial equation $s(s-1)=0$,\par
\nobreak\hfill
$u_{k=1}=a_{0}\eta\left[1+\dfrac{2-\alpha}{6}\eta^{2}
+\left(\dfrac{(2-\alpha)(12-\alpha)}{120}-\dfrac{\beta}{20}\right)\eta^{4}
+\cdots\right]$.
\solution
Expanding the differential operator gives
$(1-\eta^2)u''-2\eta u'+\alpha u+\beta\eta^2u=0$.
For $u=\sum_{j\ge0}a_j\eta^{s+j}$, the terms $-\eta^2u''$, $-2\eta u'$,
$\alpha u$ start at $\eta^s$ and $\beta\eta^2u$ at $\eta^{s+2}$, so the
two lowest powers
$\eta^{s-2},\eta^{s-1}$ come from $u''$ alone; their coefficients are
\[
s(s-1)a_0=0,\qquad s(s+1)a_1=0.
\]
Thus $s=0,1$. For $s=1$ the second equation, $2a_1=0$, forces $a_1=0$;
for $s=0$ it leaves $a_1$ free, and we take $a_1=0$ to separate the even
and odd solutions (as in 7.5.4, the odd part of the $s=0$ series is the
$s=1$ solution). At $\eta^{s+j}$, $u''$ contributes
$(s+j+2)(s+j+1)a_{j+2}$, the terms $-\eta^2u''-2\eta u'+\alpha u$
contribute $[-(s+j)(s+j-1)-2(s+j)+\alpha]a_j=[\alpha-(s+j)(s+j+1)]a_j$,
and $\beta\eta^2u$ contributes $\beta a_{j-2}$. Equating the
coefficients of $\eta^{s+j}$ to zero then gives
\[
\boxed{a_{j+2}=
\frac{[(s+j)(s+j+1)-\alpha]a_j-\beta a_{j-2}}
{(s+j+2)(s+j+1)}},\qquad a_{-2}=a_{-1}=0.
\]
For $s=1$, $j=0$ has $(s+j)(s+j+1)=2$ over the denominator
$3\cdot2=6$, and $j=2$ has $3\cdot4=12$ over $5\cdot4=20$:
\[
a_2=\frac{2-\alpha}{6}a_0,\qquad
a_4=\frac{(12-\alpha)a_2-\beta a_0}{20}
=a_0\left[\frac{(2-\alpha)(12-\alpha)}{120}-\frac\beta{20}\right].
\]
This reproduces the displayed odd solution and fulfills the request for
a series solution. For completeness, the even branch $s=0$ ($j=0$ has
$0\cdot1=0$ over $2\cdot1=2$, and $j=2$ has $2\cdot3=6$ over
$4\cdot3=12$) has
\[
a_2=-\frac\alpha2a_0,\qquad
a_4=\frac{(6-\alpha)a_2-\beta a_0}{12}
=a_0\left[\frac{\alpha(\alpha-6)}{24}-\frac\beta{12}\right].
\]
Thus a normalized pair is
\begin{align*}
u_{\rm even}&=1-\frac\alpha2\eta^2
 +\left[\frac{\alpha(\alpha-6)}{24}-\frac\beta{12}\right]\eta^4+\cdots,\\
u_{\rm odd}&=\eta\left[1+\frac{2-\alpha}{6}\eta^2
 +\left(\frac{(2-\alpha)(12-\alpha)}{120}-\frac\beta{20}\right)\eta^4+\cdots\right].
\end{align*}
The initial values $(u,u')=(1,0)$ and $(0,1)$ at zero establish
independence. Their linear combination is the general local solution;
both series converge at least for $|\eta|<1$, because a power-series
solution about an ordinary point converges in every disk about it in
which the normalized coefficients are analytic, and here
$P=-2\eta/(1-\eta^2)$ and $Q=(\alpha+\beta\eta^2)/(1-\eta^2)$ are
analytic for $|\eta|<1$ ($P$ has poles at $\eta=\pm1$).
\erratum{At special parameter values some displayed coefficients can
vanish, so these are the first three coefficients not forced to vanish
by parity, not necessarily three nonzero numbers.}

\Needspace{5\baselineskip}
\par\bigskip
\noindent\textbf{7.5.9}\quad To a good approximation, the interaction of two
nucleons may be described by a mesonic potential
\[
V=\frac{Ae^{-ax}}{x},
\]
attractive for $A$ negative. Show that the resultant Schr\"odinger wave
equation
\[
\frac{\hbar^{2}}{2m}\frac{d^{2}\psi}{dx^{2}}+(E-V)\psi=0
\]
has the following series solution \erratum{through $x^3$, allowing a
displayed coefficient to vanish at special parameter values}:
\[
\psi=a_{0}\left\{x+\frac{1}{2}A'x^{2}
+\frac{1}{6}\left[\frac{1}{2}A'^{2}-E'-aA'\right]x^{3}+\cdots\right\},
\]
where the prime indicates multiplication by $2m/\hbar^{2}$.
\solution
Put $A'=2mA/\hbar^2$, $E'=2mE/\hbar^2$. The equation becomes
\[
\psi''+\left(E'-\frac{A'e^{-ax}}x\right)\psi=0.
\]
Multiplied by $x^2$ it reads $x^2\psi''+(E'x^2-A'xe^{-ax})\psi=0$; the
coefficient of $\psi$ vanishes at $x=0$, so for $\psi\sim x^s$ only
$x^2\psi''$ contributes to the lowest power. The indicial term is
therefore $s(s-1)$, and the regular
vanishing branch has $s=1$. Write
$\psi=\sum_{j\ge0}a_jx^{j+1}$ and use
$e^{-ax}=\sum_{k\ge0}(-a)^kx^k/k!$.
Then $\psi''=\sum_{j\ge1}(j+1)ja_jx^{j-1}$ has coefficient
$(j+2)(j+1)a_{j+1}$ at $x^j$, $E'\psi$ has coefficient $E'a_{j-1}$, and
the Cauchy product
$(e^{-ax}/x)\psi=\sum_{k\ge0}\frac{(-a)^kx^k}{k!}\sum_{i\ge0}a_ix^i$
has coefficient $\sum_{k=0}^j(-a)^ka_{j-k}/k!$.
At power $x^j$, the recurrence is therefore
\[
(j+2)(j+1)a_{j+1}+E'a_{j-1}
 -A'\sum_{k=0}^j\frac{(-a)^k}{k!}a_{j-k}=0,
\qquad a_{-1}=0.
\]
For $j=0$ this gives $2a_1-A'a_0=0$, and for $j=1$ it gives
$6a_2+E'a_0-A'(a_1-aa_0)=0$, i.e.,
$6a_2=A'a_1-aA'a_0-E'a_0=(A'^2/2-aA'-E')a_0$. Consequently
\[
a_1=\frac{A'}2a_0,\qquad
a_2=\frac{a_0}{6}\left(\frac{A'^2}{2}-E'-aA'\right),
\]
and
\[
\boxed{\psi=a_0\left[x+\frac{A'}2x^2
 +\frac16\left(\frac{A'^2}{2}-E'-aA'\right)x^3+\cdots\right]}.
\]
The same coefficients follow by substituting
$e^{-ax}/x=x^{-1}-a+(a^2/2)x+\cdots$ and equating the constant and
linear terms of the differential equation. The other branch may have a
logarithm; the displayed series is the regular branch requested here.
\erratum{The qualification in the statement only makes the claim
precise: the series is determined here through $x^3$, and its displayed
coefficients need not be nonzero ($a_1=A'a_0/2$ vanishes when $A'=0$,
and $a_2$ when $E'=A'^2/2-aA'$).}

\Needspace{5\baselineskip}
\par\bigskip
\noindent\textbf{7.5.10}\quad If the parameter $b^{2}$ in Eq.~(7.53) is equal
to $2$, Eq.~(7.53) becomes
\[
y''+\frac{1}{x^{2}}y'-\frac{2}{x^{2}}y=0.
\]
From the indicial equation and the recurrence relation, derive a solution
$y=1+2x+2x^{2}$. Verify that this is indeed a solution by substituting back
into the differential equation.
\solution
Multiplying by $x^2$ gives $x^2y''+y'-2y=0$ for $x\ne0$.
For $y=\sum_{j\ge0}a_jx^{s+j}$, the lowest power $x^{s-1}$ comes only
from $y'$ (the terms $x^2y''$ and $-2y$ start at $x^s$) and has
coefficient $sa_0$. With $a_0\ne0$, the indicial equation $sa_0=0$ gives
$s=0$. Then
\[
x^2y''=\sum_{j\ge0}j(j-1)a_jx^j,\quad
y'=\sum_{j\ge0}(j+1)a_{j+1}x^j,\quad
-2y=-\sum_{j\ge0}2a_jx^j,
\]
the middle sum after replacing $j$ by $j+1$. Matching the coefficients
of $x^j$ gives
\[
[j(j-1)-2]a_j+(j+1)a_{j+1}=0,
\qquad
\boxed{a_{j+1}=\frac{2-j(j-1)}{j+1}a_j}.
\]
Thus $j=0,1,2$ give $a_1=2a_0$, $a_2=\frac22a_1=2a_0$, and
$a_3=\frac{2-2}{3}a_2=0$, so every later coefficient is zero as well.
Taking $a_0=1$ yields $\boxed{y=1+2x+2x^2}$.
The direct check is
\[
y'=2+4x,\quad y''=4,\quad
x^2y''+y'-2y=4x^2+(2+4x)-2(1+2x+2x^2)=0.
\]
The original normalized equation is undefined at zero even though this
particular solution extends there analytically. Zero is irregular
singular ($P=1/x^2$, so $xP=1/x$ is not analytic there), so finding this
terminating series does not imply that every solution admits an ordinary
Frobenius series there.

\Needspace{5\baselineskip}
\par\bigskip
\noindent\textbf{7.5.11}\quad The modified Bessel function $I_{0}(x)$
satisfies the differential equation
\[
x^{2}\frac{d^{2}}{dx^{2}}I_{0}(x)+x\frac{d}{dx}I_{0}(x)-x^{2}I_{0}(x)=0.
\]
Given that the leading term in an asymptotic expansion
\erratum{as $x\to+\infty$} is known to be
\[
I_{0}(x)\sim\frac{e^{x}}{\sqrt{2\pi x}},
\]
assume a series of the form
\[
I_{0}(x)\sim\frac{e^{x}}{\sqrt{2\pi x}}
\left\{1+b_{1}x^{-1}+b_{2}x^{-2}+\cdots\right\}.
\]
Determine the coefficients $b_{1}$ and $b_{2}$.
\par\nobreak\hfill\emph{ANS.}\enspace
$b_{1}=\dfrac{1}{8}$,\quad$b_{2}=\dfrac{9}{128}$.
\solution
Write $I_0=Ef$, where $E=e^x/\sqrt{2\pi x}$. Since
$\ln E=x-\frac12\ln x-\frac12\ln(2\pi)$, the derivatives of the
prefactor satisfy
\[
\frac{E'}E=1-\frac1{2x},\qquad
\frac{E''}E=\left(\frac{E'}E\right)'+\left(\frac{E'}E\right)^2
=1-\frac1x+\frac3{4x^2},
\]
because $(E'/E)'=1/(2x^2)$ and $(E'/E)^2=1-1/x+1/(4x^2)$.
Since $I_0'=E'f+Ef'$ and $I_0''=E''f+2E'f'+Ef''$,
substitution into $x^2I_0''+xI_0'-x^2I_0=0$ and division by $E$ give
\[
x^2f''+\left(2x^2\frac{E'}E+x\right)f'
+\left(x^2\frac{E''}E+x\frac{E'}E-x^2\right)f=0.
\]
The coefficient of $f'$ is $(2x^2-x)+x=2x^2$, and that of $f$ is
$(x^2-x+\frac34)+(x-\frac12)-x^2=\frac14$: the terms in $x^2f$ and
$xf$ cancel, leaving
\[
\boxed{x^2f''+2x^2f'+\frac14f=0}.
\]
Set $f\sim\sum_{j\ge0}b_jx^{-j}$, $b_0=1$. Termwise,
\[
x^2f''=\sum_{j\ge0}j(j+1)b_jx^{-j},\qquad
2x^2f'=-\sum_{j\ge1}2jb_jx^{1-j}=-\sum_{j\ge0}2(j+1)b_{j+1}x^{-j},
\]
the last sum after replacing $j$ by $j+1$. Since
$j(j+1)+\frac14=(j+\frac12)^2$, the coefficient of $x^{-j}$ gives
\[
[j(j+1)+\tfrac14]b_j-2(j+1)b_{j+1}=0,
\quad b_{j+1}=\frac{(j+1/2)^2}{2(j+1)}b_j.
\]
Therefore ($j=0$ and $j=1$)
\[
\boxed{b_1=\frac{1/4}{2}=\frac18},\qquad
\boxed{b_2=\frac{9/4}{4}\frac18=\frac9{128}}.
\]
These are asymptotic coefficients for large positive $x$, not a
convergent series at fixed $x$: the ratio of successive inverse-power
terms, $b_{j+1}x^{-j-1}/(b_jx^{-j})=(j+\frac12)^2/[2(j+1)x]$, behaves as
$j/(2x)$ for large $j$ and grows beyond one. Truncation
must be understood in the asymptotic sense. The qualification
$x\to+\infty$ in the statement is needed because
$I_0(x)=\sum_k(x/2)^{2k}/(k!)^2$ is an even function: as $x\to-\infty$,
$I_0(x)=I_0(|x|)\sim e^{|x|}/\sqrt{2\pi|x|}$ grows, whereas
$e^x/\sqrt{2\pi x}$ is imaginary there and tends to zero.

\Needspace{5\baselineskip}
\par\bigskip
\noindent\textbf{7.5.12}\quad The even power-series solution of Legendre's
equation is given by Exercise 8.3.1. Take $a_{0}=1$ and $n$ not an even
integer, say $n=0.5$. Calculate the partial sums of the series through
$x^{200}$, $x^{400}$, $x^{600}$, \ldots, $x^{2000}$ for $x=0.95(0.01)1.00$.
Also, write out the individual term corresponding to each of these powers.

\emph{Note.}~This calculation does not constitute proof of convergence at
$x=0.99$ or divergence at $x=1.00$, but perhaps you can see the difference in
the behavior of the sequence of partial sums for these two values of $x$.
\solution
For $n=1/2$, $n(n+1)=\frac34$ and Legendre's equation is
$(1-x^2)y''-2xy'+\tfrac34y=0$.
Insert the even series $y=\sum_{j\ge0}a_{2j}x^{2j}$ with $a_0=1$.
The coefficient of $x^{2j}$ receives $(2j+2)(2j+1)a_{2j+2}$ from $y''$
(the term $a_{2j+2}x^{2j+2}$ differentiated twice), $-2j(2j-1)a_{2j}$
from $-x^2y''$, $-4ja_{2j}$ from $-2xy'$, and $\frac34a_{2j}$ from
$\frac34y$. Since $2j(2j-1)+4j=2j(2j+1)$, it gives
\[
(2j+2)(2j+1)a_{2j+2}
+[\tfrac34-2j(2j+1)]a_{2j}=0,
\]
so
\[
\boxed{a_{2j+2}=
\frac{2j(2j+1)-3/4}{(2j+2)(2j+1)}a_{2j}}.
\]
For example $j=0$ gives $a_2=(-\frac34)/2=-3/8$, and $j=1$ gives
$a_4=\frac{6-3/4}{12}a_2=\frac7{16}\left(-\frac38\right)=-21/128$.
For each specified $x$, update the term and sum without repeatedly
forming large powers:
\[
T_0=1,\quad S_0=1,\quad
T_{j+1}=T_jx^2\frac{2j(2j+1)-3/4}{(2j+2)(2j+1)},\quad
S_{j+1}=S_j+T_{j+1}.
\]
The first table gives $S_J=\sum_{j=0}^Ja_{2j}x^{2j}$ through the power
$N=2J$. The second gives the individual term $T_J=a_Nx^N$ at that
same power. The entries are rounded to seven significant digits, a
final zero not being printed ($0.147621$ stands for $0.1476210$). The
recursion is numerically benign: ordinary double-precision arithmetic
reproduces every entry, and so does a high-precision sum of the
closed-form coefficients $a_{2j}$ given below.

\begin{center}\scriptsize
Partial sums $S_J$\par\smallskip
\setlength{\tabcolsep}{2.7pt}
\begin{tabular}{r*{6}{r}}\toprule
$N$&$x=0.95$&$0.96$&$0.97$&$0.98$&$0.99$&$1.00$\\\midrule
200&$0.2154301$&$0.1476299$&$0.0605814$&$-0.06073655$&$-0.2554509$&$-0.6500133$\\
400&$0.2154292$&$0.1476211$&$0.06048901$&$-0.06176434$&$-0.2684287$&$-0.8540902$\\
600&$0.2154292$&$0.147621$&$0.0604889$&$-0.06177416$&$-0.2694032$&$-0.9735833$\\
800&$0.2154292$&$0.147621$&$0.0604889$&$-0.06177427$&$-0.2694942$&$-1.058397$\\
1000&$0.2154292$&$0.147621$&$0.0604889$&$-0.06177428$&$-0.2695036$&$-1.124197$\\
1200&$0.2154292$&$0.147621$&$0.0604889$&$-0.06177428$&$-0.2695046$&$-1.177967$\\
1400&$0.2154292$&$0.147621$&$0.0604889$&$-0.06177428$&$-0.2695047$&$-1.223433$\\
1600&$0.2154292$&$0.147621$&$0.0604889$&$-0.06177428$&$-0.2695047$&$-1.262819$\\
1800&$0.2154292$&$0.147621$&$0.0604889$&$-0.06177428$&$-0.2695047$&$-1.297563$\\
2000&$0.2154292$&$0.147621$&$0.0604889$&$-0.06177428$&$-0.2695047$&$-1.328643$\\
\bottomrule\end{tabular}\end{center}

\begin{center}\scriptsize
Last terms $T_J$\par\smallskip
\setlength{\tabcolsep}{2.7pt}
\begin{tabular}{r*{6}{r}}\toprule
$N$&$x=0.95$&$0.96$&$0.97$&$0.98$&$0.99$&$1.00$\\\midrule
200&$-1.036298\times10^{-7}$&$-8.414151\times10^{-7}$&$-6.685141\times10^{-6}$&$-5.199707\times10^{-5}$&$-0.000396098$&$-0.002956403$\\
400&$-1.814542\times10^{-12}$&$-1.19624\times10^{-10}$&$-7.551248\times10^{-9}$&$-4.568307\times10^{-7}$&$-2.650958\times10^{-5}$&$-0.001476811$\\
600&$-4.238976\times10^{-17}$&$-2.269016\times10^{-14}$&$-1.13799\times10^{-11}$&$-5.354799\times10^{-9}$&$-2.367089\times10^{-6}$&$-0.0009842326$\\
800&$-1.114231\times10^{-21}$&$-4.842588\times10^{-18}$&$-1.92965\times10^{-14}$&$-7.062389\times10^{-11}$&$-2.378191\times10^{-7}$&$-0.000738059$\\
1000&$-3.124249\times10^{-26}$&$-1.102487\times10^{-21}$&$-3.490395\times10^{-17}$&$-9.936101\times10^{-13}$&$-2.548795\times10^{-8}$&$-0.0005903918$\\
1200&$-9.125533\times10^{-31}$&$-2.614638\times10^{-25}$&$-6.576775\times10^{-20}$&$-1.456206\times10^{-14}$&$-2.845545\times10^{-9}$&$-0.0004919624$\\
1400&$-2.741657\times10^{-35}$&$-6.378109\times10^{-29}$&$-1.274658\times10^{-22}$&$-2.195188\times10^{-16}$&$-3.267669\times10^{-10}$&$-0.0004216632$\\
1600&$-8.408677\times10^{-40}$&$-1.588298\times10^{-32}$&$-2.521936\times10^{-25}$&$-3.378161\times10^{-18}$&$-3.830633\times10^{-11}$&$-0.0003689429$\\
1800&$-2.619901\times10^{-44}$&$-4.018045\times10^{-36}$&$-5.068939\times10^{-28}$&$-5.281188\times10^{-20}$&$-4.561899\times10^{-12}$&$-0.0003279407$\\
2000&$-8.264934\times10^{-49}$&$-1.029188\times10^{-39}$&$-1.031567\times10^{-30}$&$-8.359498\times10^{-22}$&$-5.500701\times10^{-13}$&$-0.0002951405$\\
\bottomrule\end{tabular}\end{center}

The tables show the difference the Note refers to. At $x=0.99$ every
further 200 powers reduce the last term by a factor between about 8 and
15, and the partial sum is fixed to all seven digits, $-0.2695047$, from
$N=1400$ on. At $x=1.00$ the last term only halves from $N=1000$ to
$N=2000$, and the partial sums keep falling: by $0.204$ from $N=200$ to
$400$, and still by $0.031$ from $N=1800$ to $2000$. Their total fall
from $N=200$ to $2000$, $0.679$, matches
$\sum_{j=101}^{1000}0.295/j\approx0.295\ln(1000/100)=0.679$ (recall
$N=2j$), as expected from the asymptotic form $a_{2j}\approx-0.295/j$
derived below.

The convergence distinction can also be proved rather than inferred
from these tables. Since $2j(2j+1)-\frac34=4(j-\frac14)(j+\frac34)$ and
$(2j+2)(2j+1)=4(j+\frac12)(j+1)$, the recursion reads
\[
\frac{a_{2j+2}}{a_{2j}}=\frac{(j-\frac14)(j+\frac34)}{(j+\frac12)(j+1)},
\]
and because the Pochhammer symbol satisfies $(c)_{j+1}=(c)_j(c+j)$,
$(c)_0=1$, the exact coefficients are
\[
a_{2j}=\frac{(-1/4)_j(3/4)_j}{(1/2)_j j!}.
\]
The ratio of consecutive terms, $a_{2j+2}x^{2j+2}/(a_{2j}x^{2j})$,
tends to $x^2$, so the ratio test gives absolute convergence for
$|x|<1$. For $x=1$ write $(c)_j=\Gamma(c+j)/\Gamma(c)$ and use
$\Gamma(j+a)/\Gamma(j+b)\sim j^{a-b}$ as $j\to\infty$ (Stirling's formula):
\[
a_{2j}=\frac{\Gamma(1/2)}{\Gamma(-1/4)\Gamma(3/4)}
\frac{\Gamma(j-\frac14)}{\Gamma(j+\frac12)}
\frac{\Gamma(j+\frac34)}{\Gamma(j+1)}
\sim \frac{\Gamma(1/2)}{\Gamma(-1/4)\Gamma(3/4)}\,j^{-3/4}j^{-1/4}
=-\frac{\sqrt\pi}{4\Gamma(3/4)^2}\frac1j,
\]
where $\Gamma(-\frac14)=\Gamma(\frac34)/(-\frac14)=-4\Gamma(\frac34)$
by $\Gamma(z+1)=z\Gamma(z)$; the constant is $-0.29509$. (The last-term
table agrees: for $x=1$ and $N=2000$ it gives $1000\,T_{1000}=-0.2951$.)
All coefficients after $a_0$ are negative, since $a_2=-3/8$ and every
later ratio ($j\ge1$ above) is positive. By limit comparison with the
harmonic series $\sum1/j$, the series therefore diverges to $-\infty$
at $x=1$.
For every listed $x<1$, the additional factor $x^{2j}$ instead makes
the terms exponentially small. The finite numerical sequences alone
would not establish either conclusion.

\Needspace{5\baselineskip}
\par\bigskip
\noindent\textbf{7.5.13}\quad
\begin{enumerate}
\item[(a)] The odd power-series solution of Hermite's equation is given by
  Exercise 8.3.3. Take $a_{0}=1$. Evaluate this series for $\alpha=0$,
  $x=1,2,3$. Cut off your calculation after the last term calculated has
  dropped below the maximum term by a factor of $10^{6}$ or more. Set an upper
  bound to the error made in ignoring the remaining terms in the infinite
  series.
\item[(b)] As a check on the calculation of part (a), show that the Hermite
  series $y_{\mathrm{odd}}(\alpha=0)$ corresponds to
  $\erratum{\int_{0}^{x}e^{t^{2}}\,dt}$.
\item[(c)] Calculate this integral for $x=1,2,3$.
\end{enumerate}
\solution
(a) For $\alpha=0$, Hermite's equation is $y''-2xy'=0$.
The odd solution with leading coefficient $a_0=1$ has $y(0)=0$,
$y'(0)=1$. Substitute $y=\sum_{k\ge0}c_kx^{2k+1}$, $c_0=1$; after the
shift $k\to k+1$ in the second derivative,
\[
y''=\sum_{k\ge0}(2k+3)(2k+2)c_{k+1}x^{2k+1},\qquad
2xy'=\sum_{k\ge0}2(2k+1)c_kx^{2k+1},
\]
and matching the coefficients of $x^{2k+1}$ gives
\[
(2k+3)(2k+2)c_{k+1}=2(2k+1)c_k,
\qquad c_0=1.
\]
Since $2/(2k+2)=1/(k+1)$, this says
$c_{k+1}/c_k=(2k+1)/[(k+1)(2k+3)]$, which is also the ratio of
consecutive values of $1/[k!(2k+1)]$; induction therefore gives
$c_k=1/[k!(2k+1)]$. Thus
\[
\boxed{y(x)=\sum_{k=0}^\infty\frac{x^{2k+1}}{k!(2k+1)}}.
\]
For a fixed positive $x$, define $t_k=x^{2k+1}/[k!(2k+1)]$.
The iteration uses
\[
\frac{t_{k+1}}{t_k}=\frac{x^2(2k+1)}{(k+1)(2k+3)},\qquad
S_K=\sum_{k=0}^Kt_k.
\]
The ratio decreases with $k$, since apart from the constant factor $x^2$
\[
\frac{2k+1}{(k+1)(2k+3)}-\frac{2k+3}{(k+2)(2k+5)}
=\frac{4k^2+8k+1}{(k+1)(k+2)(2k+3)(2k+5)}>0,
\]
so the terms rise while the ratio exceeds one and fall afterwards. For
$x=1$ the ratio is $\frac13$ already at $k=0$; for $x=2$ it is
$\frac43,\frac65,\frac{20}{21}$ at $k=0,1,2$; for $x=3$ it is
$\frac{39}{35}>1$ at $k=6$ but $\frac{135}{136}<1$ at $k=7$. So the
largest term is $t_0=1$ for $x=1$, $t_2=3.2$ for $x=2$, and
$t_7=3^{15}/(7!\cdot15)=189.80$ for $x=3$.
Stop at the first $K$ for which $t_K\le10^{-6}\max_{j\le K}t_j$.
The thresholds are therefore $10^{-6}$, $3.2\times10^{-6}$ and
$1.898\times10^{-4}$; the terms $t_8=1.459\times10^{-6}$,
$t_{16}=1.244\times10^{-5}$ and $t_{27}=2.913\times10^{-4}$ still exceed
them, while $t_9$, $t_{17}$ and $t_{28}$ (tabulated below) do not, so
$K=9,17,28$.
For every subsequent ratio,
$t_{j+1}/t_j<x^2/(j+1)\le q:=x^2/(K+2)$ when $j\ge K+1$, the first
inequality because $(2j+1)/(2j+3)<1$. Hence $t_{K+1+i}\le q^it_{K+1}$,
and the cutoffs below have $q=\frac1{11},\frac4{19},\frac3{10}<1$, so
summing the geometric series gives the rigorous bound
\[
\boxed{0<y-S_K\le B_K:=t_{K+1}\sum_{i\ge0}q^i
=\frac{t_{K+1}}{1-x^2/(K+2)}}.
\]
This bounds the entire remaining sum, rather than assuming that the
last calculated term is itself an error bound. The tabulated bounds are
rounded upward; the sums and last terms are rounded to nearest.

\begin{center}\small
\begin{tabular}{rrrrr}\toprule
$x$&$K$&$S_K$&$t_K$&$B_K$\\\midrule
1&9&$1.4626517316055$&$1.450385\times10^{-7}$&$1.443479\times10^{-8}$\\
2&17&$16.452627042558$&$2.760027\times10^{-6}$&$7.349021\times10^{-7}$\\
3&28&$1444.5450849204$&$9.034328\times10^{-5}$&$3.869592\times10^{-5}$\\
\bottomrule\end{tabular}\end{center}

(b) \erratum{The textbook prints the integral as $\int_0^x\exp(x^2)\,dx$,
reusing $x$ as integration variable and upper limit; the dummy variable
$t$ makes its intended meaning explicit.} Since the exponential series
$e^{t^2}=\sum_{k\ge0}t^{2k}/k!$ converges uniformly
on every finite integration interval, termwise integration is valid:
\[
\int_0^xe^{t^2}dt
=\sum_{k=0}^\infty\frac1{k!}\int_0^xt^{2k}dt
=\sum_{k=0}^\infty\frac{x^{2k+1}}{k!(2k+1)}=y(x).
\]
Alternatively $y'=e^{x^2}$ directly gives $y''=2xe^{x^2}=2xy'$ and the
correct initial data $y(0)=0$, $y'(0)=1$.

(c) Numerical quadrature evaluates the integral independently of the
series: a 25-point Gauss--Legendre rule on $[0,x]$, carried out in
30-digit arithmetic, gives
\[
\begin{aligned}
\int_0^1e^{t^2}dt&=1.462651745907181608\ldots,\\
\int_0^2e^{t^2}dt&=16.45262776550723022\ldots,\\
\int_0^3e^{t^2}dt&=1444.5451228927141547\ldots.
\end{aligned}
\]
Every printed digit agrees with library values of
$(\sqrt\pi/2)\operatorname{erfi}(x)$, where
$\operatorname{erfi}(x)=-i\operatorname{erf}(ix)=(2/\sqrt\pi)\int_0^xe^{t^2}dt$
is the imaginary error function (put $u=it$ in
$\operatorname{erf}(z)=(2/\sqrt\pi)\int_0^ze^{-u^2}du$); the relative
error of the quadrature is below $10^{-26}$.
The differences from the finite sums in (a), $1.430\times10^{-8}$,
$7.229\times10^{-7}$ and $3.797\times10^{-5}$, are positive and smaller
than the bounds $B_K$, as an independent numerical check.

\Needspace{8\baselineskip}
\subsection*{7.6 Other Solutions}

\Needspace{5\baselineskip}
\par\bigskip
\noindent\textbf{7.6.1}\quad You know that the three unit vectors
$\hat{\mathbf{e}}_{x}$, $\hat{\mathbf{e}}_{y}$, and $\hat{\mathbf{e}}_{z}$ are
mutually perpendicular (orthogonal). Show that $\hat{\mathbf{e}}_{x}$,
$\hat{\mathbf{e}}_{y}$, and $\hat{\mathbf{e}}_{z}$ are linearly independent.
Specifically, show that no relation of the form of Eq.~(7.54) exists for
$\hat{\mathbf{e}}_{x}$, $\hat{\mathbf{e}}_{y}$, and $\hat{\mathbf{e}}_{z}$.
\solution
Suppose $a\hat e_x+b\hat e_y+c\hat e_z=0$.
Take the dot product successively with the three unit vectors. Their
orthonormality gives
\[
a(1)+b(0)+c(0)=0,\qquad
 a(0)+b(1)+c(0)=0,\qquad a(0)+b(0)+c(1)=0.
\]
Thus $a=b=c=0$. No nontrivial constant linear relation exists, which is
exactly linear independence. Equivalently their component matrix is the
identity and has determinant one.

\Needspace{5\baselineskip}
\par\bigskip
\noindent\textbf{7.6.2}\quad The criterion for the linear independence of
three vectors $\mathbf{A}$, $\mathbf{B}$, and $\mathbf{C}$ is that the
equation
\[
a\mathbf{A}+b\mathbf{B}+c\mathbf{C}=0,
\]
analogous to Eq.~(7.54), has no solution other than the trivial $a=b=c=0$.
Using components $\mathbf{A}=(A_{1},A_{2},A_{3})$, and so on, set up the
determinant criterion for the existence or nonexistence of a nontrivial
solution for the coefficients $a$, $b$, and $c$. Show that your criterion is
equivalent to the scalar triple product
$\mathbf{A}\cdot\mathbf{B}\times\mathbf{C}\neq0$.
\solution
The vector equation is the homogeneous component system
\[
\begin{pmatrix}A_1&B_1&C_1\\A_2&B_2&C_2\\A_3&B_3&C_3\end{pmatrix}
\begin{pmatrix}a\\b\\c\end{pmatrix}=0.
\]
It has only the zero solution precisely when its square coefficient
matrix is invertible, or equivalently when its determinant is nonzero.
Expansion down the first column gives
\begin{align*}
\det[A\ B\ C]
&=A_1(B_2C_3-B_3C_2)-A_2(B_1C_3-B_3C_1)\\
&\quad+A_3(B_1C_2-B_2C_1)
=\mathbf A\cdot(\mathbf B\times\mathbf C),
\end{align*}
because the three brackets, with the middle sign absorbed, are the
components $(B_2C_3-B_3C_2,\ B_3C_1-B_1C_3,\ B_1C_2-B_2C_1)$ of
$\mathbf B\times\mathbf C$.
Therefore
\[
\boxed{\mathbf A,\mathbf B,\mathbf C\text{ independent}
\ \Longleftrightarrow\ \mathbf A\cdot(\mathbf B\times\mathbf C)\ne0}.
\]
If the determinant is zero, the matrix has a nonzero null vector and
there is a nontrivial relation; this proves the converse as well.

\Needspace{5\baselineskip}
\par\bigskip
\noindent\textbf{7.6.3}\quad Using the Wronskian determinant, show that the
set of functions
\[
\left\{1,\frac{x^{n}}{n!},(n=1,2,\ldots,N)\right\}
\]
is linearly independent.
\solution
The set is $\{1,x,x^2/2!,\ldots,x^N/N!\}$, that is, $x^n/n!$ for
$n=0,\ldots,N$. For columns indexed by $n=0,\ldots,N$ and rows by
derivative order $j=0,\ldots,N$,
\[
\frac{d^j}{dx^j}\frac{x^n}{n!}
=\begin{cases}x^{n-j}/(n-j)!,&j\le n,\\0,&j>n.\end{cases}
\]
The Wronskian matrix is upper triangular with diagonal entries one:
\[
W=\det\begin{pmatrix}
1&x&x^2/2!&\cdots&x^N/N!\\
0&1&x&\cdots&x^{N-1}/(N-1)!\\
0&0&1&\cdots&x^{N-2}/(N-2)!\\
\vdots&\vdots&\vdots&\ddots&\vdots\\0&0&0&\cdots&1
\end{pmatrix}=\boxed{1}.
\]
If a constant linear combination of the functions were zero, its first
$N$ derivatives would be zero too. Invertibility of this matrix would
then force every coefficient to vanish.

Without the constant function, the set $\{x,x^2/2!,\ldots,x^N/N!\}$
is still independent, but its Wronskian is no longer constant. With rows
$j=0,\ldots,N-1$, the entry of column $n=1,\ldots,N$ is
$n(n-1)\cdots(n-j+1)x^{n-j}/n!$. Taking $x^n/n!$ out of column $n$ and
$x^{-j}$ out of row $j$ leaves the matrix of falling factorials
$n(n-1)\cdots(n-j+1)$. Each is $n^j$ plus lower powers of $n$ with
coefficients independent of $n$, so subtracting multiples of earlier
rows turns this matrix into the Vandermonde matrix $[n^j]$, of
determinant $\prod_{1\le i<k\le N}(k-i)=\prod_{k=2}^N(k-1)!$. Hence
\[
W\{x,x^2/2!,\ldots,x^N/N!\}
=\frac{\prod_{k=1}^{N-1}k!}{\prod_{n=1}^Nn!}\,
 x^{\sum_{n=1}^Nn-\sum_{j=0}^{N-1}j}
=\boxed{\frac{x^N}{N!}},
\]
since the two exponent sums are $N(N+1)/2$ and $N(N-1)/2$. (For $N=2$:
$x\cdot x-1\cdot x^2/2=x^2/2$.)
This Wronskian is nonzero at every $x\ne0$, enough to prove independence;
its isolated zero at zero does not imply dependence.

\Needspace{5\baselineskip}
\par\bigskip
\noindent\textbf{7.6.4}\quad If the Wronskian of two functions $y_{1}$ and
$y_{2}$ is identically zero, show by direct integration that
\[
y_{1}=cy_{2},
\]
that is, that $y_{1}$ and $y_{2}$ are linearly dependent. Assume the functions
have continuous derivatives and that at least one of the functions does not
vanish in the interval under consideration.
\erratum{Interchange the labels first, if necessary, so that $y_2$ is
nowhere zero; without this convention $y_1=cy_2$ can fail (for example
when $y_2\equiv0$ and $y_1\ne0$), although $y_1$ and $y_2$ are still
linearly dependent.}
\solution
Suppose first that $y_2$ is nowhere zero on the interval.
The quotient rule gives
\[
\left(\frac{y_1}{y_2}\right)'
=\frac{y_1'y_2-y_1y_2'}{y_2^2}=-\frac W{y_2^2}=0.
\]
Integrating from any fixed $x_*$ yields
$y_1(x)/y_2(x)=y_1(x_*)/y_2(x_*)=c$, hence
$\boxed{y_1=cy_2}$. This argument also allows $y_1$ to have zeros.

\erratum{If the nowhere-zero function is instead $y_1$, the directly
valid statement is $y_2=cy_1$. The literal requested orientation
$y_1=cy_2$ fails when $y_2\equiv0$ and $y_1\ne0$.}
In that case the same quotient calculation with roles reversed proves
linear dependence. If its constant is nonzero one can invert it and
recover the displayed orientation; if it is zero, dependence still holds
because one function is identically zero. Thus the dependence
conclusion requires no division by a function that may vanish.

\Needspace{5\baselineskip}
\par\bigskip
\noindent\textbf{7.6.5}\quad The Wronskian of two functions is found to be
zero at $x_{0}-\epsilon\le x\le x_{0}+\epsilon$ for arbitrarily small
$\epsilon>0$. Show that this Wronskian vanishes for all $x$ and that the
functions are linearly dependent.
\erratum{Assume both functions are analytic on one connected interval
containing these points, or that they solve the same normalized
second-order ODE with continuous coefficients there.}
\solution
\erratum{The textbook's statement is false for general differentiable, or even
infinitely differentiable, functions, as the counterexample below
shows. It needs analyticity on a connected
interval, or a common regular second-order ODE.}
A smooth counterexample on $\mathbb R$ is
\[
y_1(x)=1,\qquad
y_2(x)=\begin{cases}0,&x\le1,\\e^{-1/(x-1)^2},&x>1.\end{cases}
\]
To check smoothness at $1$, set $u=1/(x-1)$ on the right.
Since $du/dx=-u^2$, so that $d/dx=-u^2\,d/du$, induction shows that
every derivative has the form $p(u)e^{-u^2}$ for a polynomial $p$
(differentiating $p(u)e^{-u^2}$ gives $-u^2[p'(u)-2up(u)]e^{-u^2}$).
Each tends to zero as $u\to+\infty$, i.e.\ as $x\to1^+$; by the mean
value theorem the right-hand derivatives of every order at $1$ are then
zero, matching every derivative of the zero branch. Thus
$y_2$ is smooth.
Its Wronskian with $1$ is
\[
W=y_2'=\begin{cases}0,&x\le1,\\
\dfrac{2e^{-1/(x-1)^2}}{(x-1)^3},&x>1.\end{cases}
\]
It vanishes in every sufficiently small interval around $x_0=0$, but
not for $x>1$. The functions are independent: a relation $a+by_2=0$
first forces $a=0$ on $x\le1$, then $b=0$ on $x>1$.

Under the added analytic hypothesis, $W=y_1y_2'-y_1'y_2$ is analytic.
Its vanishing on a nonempty interval makes all its Taylor coefficients
there zero; the analytic identity theorem propagates $W=0$ through the
connected domain. If, say, $y_2$ is not identically zero, its zeros are
isolated (identity theorem again), so choose a small interval where it
is nonzero. Exercise 7.6.4 gives $y_1=cy_2$ there, and the
analytic identity theorem applied to $y_1-cy_2$ extends that relation
throughout the domain. An identically zero member gives dependence
immediately.

Alternatively, for two solutions of the same normalized ODE with
continuous $P,Q$, Abel's identity gives
$W(x)=W(x_*)e^{-\int_{x_*}^xP(t)dt}$. A single zero then makes $W$
identically zero. At any point $x_*$ the Wronskian matrix
$\bigl(\begin{smallmatrix}y_1&y_2\\y_1'&y_2'\end{smallmatrix}\bigr)$
is then singular, so it has a null vector $(a,b)\ne(0,0)$, and
$u=ay_1+by_2$ solves the ODE with $u(x_*)=u'(x_*)=0$; initial-value
uniqueness gives $u\equiv0$, which is linear dependence. The ODE or
analytic assumption is the missing reason local information can be
extended globally.

\Needspace{5\baselineskip}
\par\bigskip
\noindent\textbf{7.6.6}\quad The three functions $\sin x$, $e^{x}$, and
$e^{-x}$ are linearly independent. No one function can be written as a linear
combination of the other two. Show that the Wronskian of $\sin x$, $e^{x}$,
and $e^{-x}$ vanishes but only at isolated points.
\par\nobreak\hfill\emph{ANS.}\enspace
$W=4\sin x$,\quad$W=0$ for $x=\pm n\pi$,\quad$n=0,1,2,\ldots$.
\solution
Form the determinant with functions as columns:
\[
W=\det\begin{pmatrix}
\sin x&e^x&e^{-x}\\\cos x&e^x&-e^{-x}\\-\sin x&e^x&e^{-x}
\end{pmatrix}.
\]
Factor $e^x$ and $e^{-x}$ from the last two columns; their product is
one. Expanding the remaining determinant
\[
\det\begin{pmatrix}
\sin x&1&1\\\cos x&1&-1\\-\sin x&1&1
\end{pmatrix}
\]
along its first row, with the $2\times2$ minors $1-(-1)=2$,
$\cos x-\sin x$ and $\cos x+\sin x$, gives
\[
W=2\sin x-(\cos x-\sin x)+(\cos x+\sin x)
=\boxed{4\sin x}.
\]
It vanishes exactly at the isolated points $x=k\pi$, $k\in\mathbb Z$.
Since it is nonzero, for example, at $x=\pi/2$, the functions are
linearly independent. The special theorem that an ODE Wronskian is
either everywhere zero or nowhere zero concerns solutions of one common
regular ODE of the matching order; arbitrary independent functions can
have isolated Wronskian zeros, as this example demonstrates.

\Needspace{5\baselineskip}
\par\bigskip
\noindent\textbf{7.6.7}\quad Consider two functions $\varphi_{1}=x$ and
$\varphi_{2}=|x|$. Since $\varphi'_{1}=1$ and $\varphi'_{2}=x/|x|$,
$W(\varphi_{1},\varphi_{2})=0$ for any interval, including $[-1,+1]$. Does the
vanishing of the Wronskian over $[-1,+1]$ prove that $\varphi_{1}$ and
$\varphi_{2}$ are linearly dependent? Clearly, they are not. What is wrong?
\solution
For $x>0$, $|x|=x$ and $W=x(1)-1(x)=0$. For $x<0$, $|x|=-x$ and
$W=x(-1)-1(-x)=0$. But $|x|$ has left derivative $-1$ and right
derivative $1$ at zero, so
\[
\boxed{\text{$W(x,|x|)$ is undefined at $x=0$, not zero there.}}
\]
This is the flaw the question asks us to find; the printed argument is
an intentional trap. The continuous-derivative hypothesis needed for the
usual theorem fails.
Furthermore, both functions vanish at $x=0$, so the quotient argument of
Exercise 7.6.4, which needs one function that is nowhere zero, cannot
connect the constants on the two sides.

They are independent on that full interval: if $ax+b|x|=0$, positive
$x$ gives $a+b=0$ and negative $x$ gives $a-b=0$. Therefore
$a=b=0$. On $0<x\le1$ alone they are dependent, since $|x|=x$ there,
and on $-1\le x<0$ alone, since $|x|=-x$ there.

\Needspace{5\baselineskip}
\par\bigskip
\noindent\textbf{7.6.8}\quad Explain that linear independence does not mean
the absence of any dependence. Illustrate your argument with $\cosh x$ and
$e^{x}$.
\solution
Linear dependence requires a relation with constant coefficients,
$a\cosh x+be^x=0$, valid throughout an interval. Since
$\cosh x=(e^x+e^{-x})/2$, it would read
$(\frac a2+b)e^x+\frac a2e^{-x}=0$, and multiplication by $e^x$ gives
\[
\left(\frac a2+b\right)e^{2x}+\frac a2=0.
\]
Differentiation gives $2(a/2+b)e^{2x}=0$, hence $a/2+b=0$;
the original relation then gives $a/2=0$ and $b=0$. Thus they are
linearly independent. Their Wronskian checks this directly:
\[
W(\cosh x,e^x)=\cosh x\,e^x-\sinh x\,e^x=e^x(\cosh x-\sinh x)
=e^xe^{-x}=1.
\]
They nevertheless satisfy the nonlinear relation
\[
\boxed{2\cosh x-e^x-\frac1{e^x}=0}.
\]
Writing $Y_1=\cosh x$, $Y_2=e^x$, this is
$2Y_1-Y_2-Y_2^{-1}=0$. Linear independence excludes constant linear
relations only, not all functional or nonlinear relations.

\Needspace{5\baselineskip}
\par\bigskip
\noindent\textbf{7.6.9}\quad Legendre's differential equation
\[
(1-x^{2})y''-2xy'+n(n+1)y=0
\]
has a regular solution $P_{n}(x)$ and an irregular solution $Q_{n}(x)$. Show
that the Wronskian of $P_{n}$ and $Q_{n}$ is given by
\[
P_{n}(x)Q'_{n}(x)-P'_{n}(x)Q_{n}(x)=\frac{A_{n}}{1-x^{2}},
\]
with $A_{n}$ independent of $x$.
\solution
Let $u=P_n$, $v=Q_n$ and $W=uv'-u'v$ on an interval not crossing
$x=\pm1$. Differentiate so the cross terms cancel:
\[
W'=u'v'+uv''-u''v-u'v'=uv''-u''v.
\]
The normalized Legendre equation gives
$y''=2xy'/(1-x^2)-n(n+1)y/(1-x^2)$. Substitution for $u''$ and $v''$
yields
\[
W'=\frac{2x(uv'-u'v)-n(n+1)(uv-uv)}{1-x^2}=\frac{2x}{1-x^2}W.
\]
Equivalently,
$[(1-x^2)W]'=-2xW+(1-x^2)W'=0$. Integrating gives $(1-x^2)W=A_n$, a
constant, that is,
\[
\boxed{W(P_n,Q_n)=\frac{A_n}{1-x^2}},
\]
with $A_n$ constant on the interval. Its value depends on the
normalization of the chosen solutions; the differential equation alone
determines the $x$ dependence.

\Needspace{5\baselineskip}
\par\bigskip
\noindent\textbf{7.6.10}\quad Show, by means of the Wronskian, that a linear,
second-order, homogeneous ODE of the form
\[
y''(x)+P(x)y'(x)+Q(x)y(x)=0
\]
cannot have three independent solutions.
\erratum{Assume $P,Q$ are continuous on a connected interval.}

\emph{Hint.}~Assume a third solution and show that the Wronskian vanishes for
all $x$.
\solution
\erratum{Work on a connected interval where $P,Q$ are continuous, so the
normalized second-order initial-value problem is unique.}
Suppose $y_1,y_2,y_3$ solve the equation. In their $3\times3$ Wronskian,
\[
W_3=\det\begin{pmatrix}y_1&y_2&y_3\\y_1'&y_2'&y_3'\\y_1''&y_2''&y_3''\end{pmatrix},
\]
the third row equals $-Q(x)$ times the first row minus $P(x)$ times
the second. Hence $W_3(x)=0$ for every $x$.

A vanishing Wronskian alone does not prove dependence of arbitrary
differentiable functions: $x^2$ and $x|x|$, whose derivatives are $2x$
and $2|x|$, have $W=x^2\cdot2|x|-2x\cdot x|x|=0$ everywhere, yet they are
independent on $[-1,1]$ by the argument of Exercise 7.6.7. So complete
the argument with the ODE: choose any point $x_0$ of the interval. The
two homogeneous equations
\[
\sum_{j=1}^3c_jy_j(x_0)=0,\qquad
\sum_{j=1}^3c_jy_j'(x_0)=0
\]
have a nonzero constant solution $(c_1,c_2,c_3)$ because there are
three unknowns and only two equations. The function
$u=\sum c_jy_j$ solves the same ODE and has $u(x_0)=u'(x_0)=0$.
Uniqueness forces $u\equiv0$. This supplies a nontrivial constant
linear relation, proving that three independent solutions cannot exist.

\Needspace{5\baselineskip}
\par\bigskip
\noindent\textbf{7.6.11}\quad Show the following when the linear second-order
differential equation $py''+qy'+ry=0$ is expressed in self-adjoint form:
\erratum{Use the coefficients after conversion, so $q=p'$, and work on
an interval with $p\ne0$; part (b) initially requires $y_1\ne0$.}
\begin{enumerate}
\item[(a)] The Wronskian is equal to a constant divided by $p$:
  \[
  W(x)=\frac{C}{p(x)}.
  \]
\item[(b)] A second solution $y_{2}(x)$ is obtained from a first solution
  $y_{1}(x)$ as
  \[
  y_{2}(x)=Cy_{1}(x)\int^{x}\frac{dt}{p(t)[y_{1}(t)]^{2}}.
  \]
\end{enumerate}
\solution
In self-adjoint form the equation is
\[
(py')'+ry=0,\qquad\text{so }py''+p'y'+ry=0.
\]
Here $p$ denotes the coefficient after any integrating factor has been
applied; the self-adjoint condition is $q=p'$. Work where $p\ne0$.

(a) For solutions $y_1,y_2$, subtract $y_2$ times the first equation
from $y_1$ times the second. The terms involving $r$ cancel:
\[
p(y_1y_2''-y_1''y_2)+p'(y_1y_2'-y_1'y_2)=0.
\]
Since $W=y_1y_2'-y_1'y_2$ has $W'=y_1y_2''-y_1''y_2$ (the terms
$y_1'y_2'$ cancel), this is $pW'+p'W=(pW)'=0$, and
\[
\boxed{W=C/p}.
\]
(b) On a subinterval where $y_1\ne0$,
\[
\left(\frac{y_2}{y_1}\right)'
=\frac{y_1y_2'-y_1'y_2}{y_1^2}=\frac{C}{p y_1^2}.
\]
Integration gives
\[
\boxed{y_2=y_1\left[D+C\int^x\frac{dt}{p(t)y_1(t)^2}\right]}.
\]
The $D$ term only repeats the known solution; set $D=0$ and choose
$C\ne0$ for an independent second solution. The formula is initially
local around points with $y_1\ne0$ and may then be continued by the ODE.

\Needspace{5\baselineskip}
\par\bigskip
\noindent\textbf{7.6.12}\quad Transform our linear, second-order ODE
\[
y''+P(x)y'+Q(x)y=0
\]
\erratum{with $P$ continuously differentiable and $Q$ continuous on the
interval}, by the substitution
\[
y=z\exp\left[-\frac{1}{2}\int^{x}P(t)\,dt\right]
\]
and show that the resulting differential equation for $z$ is
\[
z''+q(x)z=0,
\]
where
\[
q(x)=Q(x)-\frac{1}{2}P'(x)-\frac{1}{4}P^{2}(x).
\]
\emph{Note.}~This substitution can be derived by the technique of
Exercise 7.6.25.
\solution
Set $E(x)=\exp[-\tfrac12\int^xP(t)dt]$, so $y=Ez$ and, differentiating
$E'=-PE/2$ once more,
\[
E'=-\frac P2E,\qquad E''=-\frac{P'}2E-\frac P2E'
=\left(-\frac{P'}2+\frac{P^2}4\right)E.
\]
The product rule gives
\[
y'=E\left(z'-\frac P2z\right),\qquad
y''=E\left[z''-Pz'+\left(-\frac{P'}2+\frac{P^2}4\right)z\right].
\]
Substitute all three terms of the ODE:
\begin{align*}
y''+Py'+Qy
&=E\left[z''-Pz'+\left(-\frac{P'}2+\frac{P^2}4\right)z
 +Pz'-\frac{P^2}2z+Qz\right]\\
&=E\left[z''+\left(Q-\frac{P'}2-\frac{P^2}4\right)z\right].
\end{align*}
Since $E$ never vanishes, division by it proves
\[
\boxed{z''+qz=0,\qquad q=Q-\tfrac12P'-\tfrac14P^2}.
\]
This calculation requires $P$ differentiable where the substitution is
used, in addition to the usual existence assumptions on the coefficients.

\Needspace{5\baselineskip}
\par\bigskip
\noindent\textbf{7.6.13}\quad Use the result of Exercise 7.6.12 to show that
the replacement of $\varphi(r)$ by $r\varphi(r)$ may be expected to eliminate
the first derivative from the Laplacian in spherical polar coordinates. See
also Exercise 3.10.34.
\solution
The radial part of the spherical Laplacian is
\[
\frac1{r^2}\frac d{dr}\left(r^2\frac{d\varphi}{dr}\right)
=\varphi''+\frac2r\varphi',\qquad r>0.
\]
Thus $P(r)=2/r$ and the factor in 7.6.12 is
$\exp[-\tfrac12\int(2/r)dr]=e^{-\ln r}=1/r$, up to a constant.
The new dependent variable is therefore $u=r\varphi$, or
$\varphi=u/r$. Direct differentiation confirms the cancellation:
\[
\varphi'=\frac{u'}r-\frac u{r^2},\qquad
\varphi''=\frac{u''}r-\frac{2u'}{r^2}+\frac{2u}{r^3},
\]
so that
\[
\varphi''+\frac2r\varphi'
=\frac{u''}r-\frac{2u'}{r^2}+\frac{2u}{r^3}+\frac{2u'}{r^2}-\frac{2u}{r^3},
\qquad
\boxed{\varphi''+\frac2r\varphi'=\frac{u''}r}.
\]
The potential correction from the substitution also vanishes:
$-P'/2-P^2/4=1/r^2-1/r^2=0$, so $q=Q$ in 7.6.12.
If angular dependence is present, it is unaffected: writing
$\nabla^2=r^{-2}\partial_r(r^2\partial_r)+r^{-2}\Delta_\Omega$, where
$\Delta_\Omega$ is the angular operator in $\theta$ and the azimuth
(it does not act on $r$),
$\nabla^2(u/r)=u_{rr}/r+\Delta_\Omega u/r^3$.
The eliminated derivative is the first \emph{radial} derivative. At the
origin a separate condition is needed: for $\varphi=u/r$ to stay finite
there, $u=r\varphi$ must vanish at $r=0$.

\Needspace{5\baselineskip}
\par\bigskip
\noindent\textbf{7.6.14}\quad By direct differentiation and substitution show
that
\[
y_{2}(x)=y_{1}(x)\int^{x}\frac{\exp\left[-\int^{s}P(t)\,dt\right]}
{[y_{1}(s)]^{2}}\,ds
\]
satisfies, like $y_{1}(x)$, the ODE
\[
y''_{2}(x)+P(x)y'_{2}(x)+Q(x)y_{2}(x)=0.
\]
\emph{Note.}~The Leibniz formula for the derivative of an integral is
\[
\frac{d}{d\alpha}\int_{g(\alpha)}^{h(\alpha)}f(x,\alpha)\,dx
=\int_{g(\alpha)}^{h(\alpha)}\frac{\partial f(x,\alpha)}{\partial\alpha}\,dx
+f[h(\alpha),\alpha]\frac{dh(\alpha)}{d\alpha}
-f[g(\alpha),\alpha]\frac{dg(\alpha)}{d\alpha}.
\]
\erratum{Apply the integral formula on an interval with continuous $P,Q$
and a nowhere-zero first solution $y_1$; choose all integration limits
within that interval.}
\solution
On an interval where $y_1\ne0$, abbreviate
\[
E=e^{-\int^xP(t)dt},\qquad h=E/y_1^2,\qquad
J=\int^xh(s)ds,
\]
so $E'=-PE$, $J'=h$ (the Leibniz formula of the note, with only the
upper limit depending on $x$), and $y_2=y_1J$. Differentiate
explicitly; $h'=E'/y_1^2-2Ey_1'/y_1^3$ and the product rule give
\[
h'=-Ph-2\frac{y_1'}{y_1}h,
\quad y_2'=y_1'J+y_1h,
\quad y_2''=y_1''J+2y_1'h+y_1h'.
\]
Substitution groups the terms as
\begin{align*}
y_2''+Py_2'+Qy_2
&=J(y_1''+Py_1'+Qy_1)+(2y_1'+Py_1)h+y_1h'\\
&=0+(2y_1'+Py_1)h
  +y_1\left(-P-2\frac{y_1'}{y_1}\right)h=0.
\end{align*}
Thus the displayed expression is a solution. Its Wronskian with $y_1$
is $y_1(y_1'J+y_1h)-y_1'(y_1J)=y_1^2h=E\ne0$, independently
confirming that it is a second, independent solution.

\Needspace{5\baselineskip}
\par\bigskip
\noindent\textbf{7.6.15}\quad In the equation
\[
y_{2}(x)=y_{1}(x)\int^{x}\frac{\exp\left[-\int^{s}P(t)\,dt\right]}
{[y_{1}(s)]^{2}}\,ds,
\]
$y_{1}(x)$ satisfies
\[
y''_{1}+P(x)y'_{1}+Q(x)y_{1}=0.
\]
The function $y_{2}(x)$ is a linearly independent second solution of the same
equation. Show that the inclusion of lower limits on the two integrals leads
to nothing new, that is, that it generates only an overall constant factor and
a constant multiple of the known solution $y_{1}(x)$.
\erratum{Apply the integral formula on an interval with continuous $P,Q$
and a nowhere-zero first solution $y_1$; choose all integration limits
within that interval.}
\solution
Distinguish the inner and outer lower limits by defining
\[
E_a(s)=\exp\left[-\int_a^sP(t)dt\right],\qquad
Y_{a,c}(x)=y_1(x)\int_c^x\frac{E_a(s)}{y_1(s)^2}ds.
\]
For another inner lower limit $b$,
\[
\int_a^sP=\int_b^sP-\int_b^aP,\qquad
E_a(s)=K E_b(s),\quad K=\exp\left(\int_b^aP(t)dt\right)\ne0.
\]
Changing the outer lower limit from $c$ to $d$, with
$\int_c^x=\int_d^x-\int_d^c$, then gives
\begin{align*}
Y_{a,c}
&=K y_1\int_c^x\frac{E_b(s)}{y_1(s)^2}ds\\
&=K y_1\left[\int_d^x\frac{E_b(s)}{y_1(s)^2}ds
              -\int_d^c\frac{E_b(s)}{y_1(s)^2}ds\right]\\
&=\boxed{K Y_{b,d}+D y_1},\qquad
D=-K\int_d^c\frac{E_b(s)}{y_1(s)^2}ds.
\end{align*}
Both $K$ and $D$ are constants. Hence lower limits only rescale the
second solution and add a multiple of the first. The limits must lie in
an interval where the integrals and the nonzero-denominator formula are
well defined.

\Needspace{5\baselineskip}
\par\bigskip
\noindent\textbf{7.6.16}\quad Given that one solution of
\[
R''+\frac{1}{r}R'-\frac{m^{2}}{r^{2}}R=0
\]
is $R=r^{m}$, show that Eq.~(7.67) predicts a second solution,
$R=r^{-m}$ \erratum{for $m\ne0$, and $R=\ln r$ for $m=0$, on $r>0$}.
\solution
On $r>0$, $P=1/r$ and $e^{-\int Pdr}=1/r$, up to a constant.
With $R_1=r^m$, reduction of order gives
\[
R_2=r^m\int^r\frac{ds/s}{s^{2m}}
=r^m\int^rs^{-2m-1}ds.
\]
For $m\ne0$, $\int^rs^{-2m-1}ds=-r^{-2m}/(2m)$, so
\[
R_2=-\frac1{2m}r^{-m},
\]
and rescaling yields the requested independent solution
$\boxed{R_2=r^{-m}}$. With $R_2'=-mr^{-m-1}$ and
$R_2''=m(m+1)r^{-m-2}$, direct substitution gives
$m(m+1)r^{-m-2}-mr^{-m-2}-m^2r^{-m-2}=0$.
\erratum{When $m=0$, $r^m=r^{-m}=1$ does not provide two independent
solutions. The integral becomes $\int ds/s$, and the second solution
is $\ln r$.} Its derivatives $1/r,-1/r^2$ satisfy the $m=0$ equation.

\Needspace{5\baselineskip}
\par\bigskip
\noindent\textbf{7.6.17}\quad Using
\[
y_{1}(x)=\sum_{n=0}^{\infty}\frac{(-1)^{n}}{(2n+1)!}x^{2n+1}
\]
as a solution of the linear oscillator equation, follow the analysis that
proceeds through Eq.~(7.81) and show that in that equation $c_{n}=0$, so that
in this case the second solution does not contain a logarithmic term.
\solution
The given first solution is $y_1=\sin x=xA(x^2)$ with
$A(w)=\sum_{m\ge0}(-1)^mw^m/(2m+1)!$, so $A(0)=1$; the oscillator
equation $y''+y=0$ has $P=0$, so $e^{-\int P\,dx}=1$.
Its reduction-of-order integrand is
\[
\frac1{y_1^2}=\frac1{x^2}\frac1{A(x^2)^2}
=\frac1{x^2}\sum_{j\ge0}b_jx^{2j}
=\sum_{j\ge0}b_jx^{2j-2}.
\]
The reciprocal series exists near zero because $A(0)\ne0$.
A logarithm can arise on termwise integration only from a power
$x^{-1}$. But $2j-2=-1$ would require $j=1/2$, not an integer.
Thus the logarithm coefficient denoted $c_n$ in the referenced analysis
is zero: the indicial roots of $y''+y=0$ at $x=0$ are $1$ and $0$, so
$n=1$, the integrand is $x^{-2}(1+c_1x+c_2x^2+\cdots)$, and $c_1$, the
coefficient of $x$ in the even series $1/A(x^2)^2$, vanishes. Indeed,
\[
y_2=\sin x\int^x\frac{ds}{\sin^2s}
=\sin x(-\cot x)=-\cos x,
\]
since $(\cot s)'=-1/\sin^2s$; the integration constant adds a multiple
of $\sin x$. The Taylor series
$-1+x^2/2!-x^4/4!+\cdots$ contains no logarithm and verifies the
coefficient argument directly.

\Needspace{5\baselineskip}
\par\bigskip
\noindent\textbf{7.6.18}\quad Show that when $n$ is not an integer in Bessel's
ODE, Eq.~(7.40), the second solution of Bessel's equation, obtained from
Eq.~(7.67), does not contain a logarithmic term.
\solution
For a real noninteger order set $\nu=|n|$; for a complex noninteger
order the same proof works with $\nu=n$. Bessel's equation
$x^2y''+xy'+(x^2-n^2)y=0$ depends only on $n^2=\nu^2$. Substituting
$y=\sum_{j\ge0}a_jx^{\nu+j}$ gives $[(\nu+j)^2-\nu^2]a_j+a_{j-2}=0$,
that is, $j(2\nu+j)a_j=-a_{j-2}$, with $a_{-1}=a_{-2}=0$. Here
$2\nu+j\ne0$ for $j\ge1$, so $a_1=0$ and every odd coefficient
vanishes: the Frobenius solution with exponent $\nu$ has the form
\[
y_1=x^\nu A(x^2),\qquad A(0)=a_0\ne0.
\]
Since $P=1/x$ and $e^{-\int P\,dx}=1/x$,
\[
\frac{e^{-\int Pdx}}{y_1^2}
=x^{-2\nu-1}\frac1{A(x^2)^2}
=\sum_{k\ge0}b_kx^{2k-2\nu-1}.
\]
A logarithm could arise only if $2k-2\nu-1=-1$, or $k=\nu$.
There is no such nonnegative integer $k$ when $\nu$ is noninteger.
Termwise integration therefore gives
\[
y_2=x^\nu A(x^2)\sum_{k\ge0}
\frac{b_k}{2k-2\nu}x^{2k-2\nu}
=x^{-\nu}\times\text{a power series in }x^2,
\]
with no logarithm. The power-series integrations are local, and a
consistent branch of $x^\nu$ is used for complex $x$.
The leading coefficient is nonzero, so this is
independent of $x^\nu A(x^2)$ and proportional, up to addition of that
first solution, to $J_{-\nu}$. This also explains why half-integer
orders, though their indicial roots differ by an integer, still produce
no logarithm: parity prevents the resonant $x^{-1}$ term.

\Needspace{5\baselineskip}
\par\bigskip
\noindent\textbf{7.6.19}\quad
\begin{enumerate}
\item[(a)] One solution of Hermite's differential equation
  \[
  y''-2xy'+2\alpha y=0
  \]
  for $\alpha=0$ is $y_{1}(x)=1$. Find a second solution, $y_{2}(x)$, using
  Eq.~(7.67). Show that your second solution is equivalent to
  $y_{\mathrm{odd}}$ (Exercise 8.3.3).
\item[(b)] Find a second solution for $\alpha=1$, where $y_{1}(x)=x$, using
  Eq.~(7.67). Show that your second solution is equivalent to
  $y_{\mathrm{even}}$ (Exercise 8.3.3).
\end{enumerate}
\solution
For Hermite's equation $P=-2x$, so $e^{-\int Pdx}=e^{x^2}$.

(a) When $\alpha=0$ and $y_1=1$,
\[
\boxed{y_2=\int_0^xe^{t^2}dt
=\sum_{k=0}^\infty\frac{x^{2k+1}}{k!(2k+1)}}.
\]
The exponential series converges uniformly on finite intervals,
justifying its termwise integral. This is the normalized odd solution:
$y_2(0)=0$, $y_2'(0)=1$, and $y_2''-2xy_2'=0$ directly. To compare
with Exercise 8.3.3, substitute $y=\sum_ja_jx^j$ into Hermite's
equation: $(j+2)(j+1)a_{j+2}-2ja_j+2\alpha a_j=0$, so
\[
a_{j+2}=\frac{2(j-\alpha)}{(j+2)(j+1)}a_j.
\]
For $\alpha=0$ and $j=2k+1$ this gives
$a_{2k+3}/a_{2k+1}=(2k+1)/[(k+1)(2k+3)]$, the ratio of consecutive
coefficients $1/[k!(2k+1)]$ above (as in Exercise 7.5.13); with the
same leading coefficient the two odd series coincide:
$x+x^3/3+x^5/10+\cdots$.

(b) With $\alpha=1$ and $y_1=x$, reduction of order gives
$x\int^xe^{t^2}/t^2\,dt$. Differentiate $e^{t^2}/t$:
\[
\frac d{dt}\left(\frac{e^{t^2}}t\right)
=2e^{t^2}-\frac{e^{t^2}}{t^2},
\quad
\int^x\frac{e^{t^2}}{t^2}dt=-\frac{e^{x^2}}x+2\int_0^xe^{t^2}dt.
\]
Multiplying by $x$ gives a solution with value $-1$ at zero. Reverse
its sign to normalize the even solution to one:
\[
\boxed{y_{\rm even}=e^{x^2}-2x\int_0^xe^{t^2}dt}.
\]
Its series follows from $e^{x^2}=\sum_kx^{2k}/k!$ and the series of
part (a), with $k\to k-1$ in the second sum:
\begin{align*}
y_{\rm even}
&=\sum_{k\ge0}\frac{x^{2k}}{k!}
 -2\sum_{k\ge0}\frac{x^{2k+2}}{k!(2k+1)}\\
&=1+\sum_{k\ge1}\left[\frac1{k!}
 -\frac2{(k-1)!(2k-1)}\right]x^{2k}\\
&=1-\sum_{k\ge1}\frac{x^{2k}}{k!(2k-1)}
=1-x^2-\frac{x^4}{6}-\frac{x^6}{30}-\cdots,
\end{align*}
because $2/(k-1)!=2k/k!$ and
\[
\frac1{k!}-\frac{2k}{k!(2k-1)}=\frac{(2k-1)-2k}{k!(2k-1)}
=-\frac1{k!(2k-1)}.
\]
This is the even Hermite series at $\alpha=1$: the recursion of part (a)
with $\alpha=1$, $j=2k$, gives $a_2=-a_0$ and, for $k\ge1$,
$a_{2k+2}/a_{2k}=(2k-1)/[(k+1)(2k+1)]$, which is exactly the ratio of
consecutive coefficients $-1/[k!(2k-1)]$. Setting
$J=\int_0^xe^{t^2}dt$, so $J'=e^{x^2}$, gives
$y'=2xe^{x^2}-2J-2xe^{x^2}=-2J$, $y''=-2e^{x^2}$, and
$y''-2xy'+2y=-2e^{x^2}+4xJ+2(e^{x^2}-2xJ)=0$.
The integral construction is initially used on $x\ne0$; the displayed
even function extends smoothly through zero.

\Needspace{5\baselineskip}
\par\bigskip
\noindent\textbf{7.6.20}\quad One solution of Laguerre's differential equation
\[
xy''+(1-x)y'+ny=0
\]
for $n=0$ is $y_{1}(x)=1$. Using Eq.~(7.67), develop a second, linearly
independent solution. Exhibit the logarithmic term explicitly.
\solution
Normalize the $n=0$ equation to $y''+(1/x-1)y'=0$ on $x\ne0$.
Since $\int(1/x-1)dx=\ln|x|-x$,
\[
e^{-\int(1/x-1)dx}=\frac{e^x}{|x|}=\pm\frac{e^x}{x}
\]
on either real half-line; the constant sign is dropped. With $y_1=1$,
$y_1^2=1$, Eq.~(7.67) gives
\[
y_2=\int^x\frac{e^s}{s}ds.
\]
Expand $e^s/s=s^{-1}+\sum_{k\ge1}s^{k-1}/k!$ and integrate:
\[
\boxed{y_2=\ln|x|+\sum_{k=1}^\infty\frac{x^k}{k\,k!}}.
\]
An additive constant repeats $y_1$ and is omitted. The series factor
converges for all finite $x$, by its factorial denominator.
Its derivative is $e^x/x$, so $W(1,y_2)=e^x/x\ne0$, proving
independence. The logarithmic singularity at zero is explicit.

\Needspace{5\baselineskip}
\par\bigskip
\noindent\textbf{7.6.21}\quad For Laguerre's equation with $n=0$,
\[
y_{2}(x)=\int^{x}\frac{e^{s}}{s}\,ds.
\]
\begin{enumerate}
\item[(a)] Write $y_{2}(x)$ as a logarithm plus a power series.
\item[(b)] Verify that the integral form of $y_{2}(x)$, previously given, is a
  solution of Laguerre's equation ($n=0$) by direct differentiation of the
  integral and substitution into the differential equation.
\item[(c)] Verify that the series form of $y_{2}(x)$, part (a), is a solution
  by differentiating the series and substituting back into Laguerre's
  equation.
\end{enumerate}
\solution
(a) Separate the integrand's pole before integrating:
\[
\frac{e^s}s=\frac1s+\sum_{k=1}^\infty\frac{s^{k-1}}{k!},
\qquad
\boxed{y_2=\ln|x|+\sum_{k=1}^\infty\frac{x^k}{k\,k!}+C},\quad x\ne0.
\]
The factorial series and its derivatives converge locally uniformly,
so these integrations and differentiations are justified away from zero.

(b) The fundamental theorem gives $y_2'=e^x/x$, and the quotient rule
gives $y_2''=e^x/x-e^x/x^2$. Consequently
\[
xy_2''+(1-x)y_2'
=e^x-\frac{e^x}x+\frac{e^x}x-e^x=0.
\]
(c) Differentiate the series explicitly, using
$[x^k/(k\,k!)]'=x^{k-1}/k!$:
\[
y_2'=\frac1x+\sum_{k\ge1}\frac{x^{k-1}}{k!},\qquad
y_2''=-\frac1{x^2}+\sum_{k\ge1}\frac{(k-1)x^{k-2}}{k!}.
\]
Hence
\[
xy_2''=-\frac1x+\sum_{k\ge1}\frac{(k-1)x^{k-1}}{k!},\qquad
(1-x)y_2'=\frac1x-1+\sum_{k\ge1}\frac{x^{k-1}}{k!}
 -\sum_{k\ge1}\frac{x^k}{k!}.
\]
Adding, the $1/x$ terms cancel and $(k-1)+1=k$, so
\begin{align*}
xy_2''+(1-x)y_2'
&=-1+\sum_{k\ge1}\frac{kx^{k-1}}{k!}
      -\sum_{k\ge1}\frac{x^k}{k!}\\
&=-1+\sum_{k\ge1}\frac{x^{k-1}}{(k-1)!}-(e^x-1)
=-1+e^x-(e^x-1)=0.
\end{align*}
This verifies the series independently of recognizing it as the integral.

\Needspace{5\baselineskip}
\par\bigskip
\noindent\textbf{7.6.22}\quad One solution of the Chebyshev equation
\[
(1-x^{2})y''-xy'+n^{2}y=0
\]
for $n=0$ is $y_{1}=1$.
\begin{enumerate}
\item[(a)] Using Eq.~(7.67), develop a second, linearly independent solution.
\item[(b)] Find a second solution by direct integration of the Chebyshev
  equation.
\end{enumerate}
\emph{Hint.}~Let $v=y'$ and integrate. Compare your result with the second
solution given in Section 18.4.
\par\nobreak\hfill\emph{ANS.}\enspace
(a) $y_{2}=\sin^{-1}x$ \erratum{on $-1<x<1$}.\quad
(b) The second solution, $V_n(x)$, is not defined for $n=0$
\erratum{under the nonzero-$n$ convention of Section 18.4}.
\solution
Work on $-1<x<1$. (a) The normalized coefficient is
$P=-x/(1-x^2)$; with $w=1-x^2$, $dw=-2x\,dx$,
\[
\int Pdx=\int\frac{dw}{2w}=\tfrac12\ln(1-x^2),\qquad
e^{-\int Pdx}=(1-x^2)^{-1/2}.
\]
Since $y_1=1$, reduction of order yields
\[
\boxed{y_2=\int_0^x\frac{ds}{\sqrt{1-s^2}}=\arcsin x}.
\]
Its Wronskian with $1$ is $1/\sqrt{1-x^2}\ne0$.

(b) Put $v=y'$ in the $n=0$ equation:
\[
(1-x^2)v'-xv=0,\quad
\frac{dv}{v}=\frac{x\,dx}{1-x^2},\quad
\ln|v|=-\tfrac12\ln(1-x^2)+C.
\]
Therefore $v=A/\sqrt{1-x^2}$ and
$\boxed{y=A\arcsin x+B}$, reproducing the same pair.

Section 18.4 explicitly introduces $V_n(x)=\sin(n\arccos x)$ for
$n\ne0$, so the printed answer follows its convention. The sine expression
has a continuous extension to $n=0$, but this gives the zero function,
which cannot be an independent second solution. The expression
$V_n=\sqrt{1-x^2}\,U_{n-1}$ also falls outside the usual nonnegative
polynomial indexing when $n=0$.
A nonzero limiting second solution is
$\lim_{n\to0}V_n(x)/n=\arccos x=\pi/2-\arcsin x$, differing only
by a sign and a multiple of the first solution.

\Needspace{5\baselineskip}
\par\bigskip
\noindent\textbf{7.6.23}\quad One solution of the Chebyshev equation
\[
(1-x^{2})y''-xy'+n^{2}y=0
\]
for $n=1$ is $y_{1}(x)=x$. Set up the Wronskian double integral solution and
derive a second solution, $y_{2}(x)$.
\par\nobreak\hfill\emph{ANS.}\enspace $y_{2}=-(1-x^{2})^{1/2}$.
\solution
On $-1<x<1$, use $P=-x/(1-x^2)$ and $y_1=x$. As in Exercise 7.6.22,
the inner integral is $\int^sP(t)dt=\frac12\ln(1-s^2)$, so
$\exp[-\int^sP(t)dt]=(1-s^2)^{-1/2}$, and the
Wronskian double integral gives, initially on either side of zero,
\[
y_2=x\int^x\frac{\exp[-\int^s(-t/(1-t^2))dt]}{s^2}ds
=x\int^x\frac{ds}{s^2\sqrt{1-s^2}}.
\]
The antiderivative follows directly from the product rule, since
$(\sqrt{1-s^2})'=-s/\sqrt{1-s^2}$:
\[
\frac d{ds}\left(-\frac{\sqrt{1-s^2}}s\right)
=\frac1{\sqrt{1-s^2}}+\frac{\sqrt{1-s^2}}{s^2}
=\frac{s^2+(1-s^2)}{s^2\sqrt{1-s^2}}
=\frac1{s^2\sqrt{1-s^2}}.
\]
Thus $y_2=x\cdot(-\sqrt{1-x^2}/x)$, that is,
\[
\boxed{y_2=-\sqrt{1-x^2}},
\]
up to a multiple of $x$ from the integration constant. Its derivatives
are, by the chain rule and then the quotient rule,
\[
y_2'=\frac x{\sqrt{1-x^2}},\qquad
y_2''=\frac{\sqrt{1-x^2}+x^2/\sqrt{1-x^2}}{1-x^2}=(1-x^2)^{-3/2},
\]
so substitution gives
\[
(1-x^2)y_2''-xy_2'+y_2
=\frac{1-x^2}{\sqrt{1-x^2}}-\sqrt{1-x^2}=0.
\]
Finally $W(x,y_2)=xy_2'-y_2=x^2/\sqrt{1-x^2}+\sqrt{1-x^2}
=1/\sqrt{1-x^2}\ne0$. Although the reduction
integrand was singular at zero, the final solution and this Wronskian
extend there regularly.

\Needspace{5\baselineskip}
\par\bigskip
\noindent\textbf{7.6.24}\quad The radial Schr\"odinger wave equation for a
spherically symmetric potential can be written in the form
\[
\left[-\frac{\hbar^{2}}{2m}\frac{d^{2}}{dr^{2}}
+l(l+1)\frac{\hbar^{2}}{2mr^{2}}+V(r)\right]y(r)=Ey(r).
\]
\erratum{Take $l=0,1,2,\ldots$ and $r>0$.}
The potential energy $V(r)$ may be expanded about the origin as
\[
V(r)=\frac{b_{-1}}{r}+b_{0}+b_{1}r+\cdots.
\]
\begin{enumerate}
\item[(a)] Show that there is one (regular) solution $y_{1}(r)$ starting with
  $r^{l+1}$.
\item[(b)] From Eq.~(7.69) show that the irregular solution $y_{2}(r)$
  \erratum{has leading behavior $r^{-l}$, diverging for $l\ge1$ but
approaching a nonzero constant for $l=0$.}
\end{enumerate}
\solution
Let $\kappa=2m/\hbar^2$ and take the physical $l=0,1,2,\ldots$.
Multiplying by $-\kappa$ and collecting all terms on the left, the
equation becomes
\[
y''-\frac{l(l+1)}{r^2}y+\kappa[E-V(r)]y=0.
\]
(a) Insert a Frobenius trial $y=\sum_{j\ge0}a_jr^{s+j}$, $a_0\ne0$.
The potential has at most a simple pole, so $\kappa(E-V)y$ begins at
$r^{s-1}$, and the lowest power $r^{s-2}$ comes only from
$y''-l(l+1)y/r^2$, with coefficient
\[
[s(s-1)-l(l+1)]a_0=(s-l-1)(s+l)a_0=0.
\]
The roots are $l+1$ and $-l$. For $s=l+1$, the coefficient of $a_j$
in the recurrence is
\[
(s+j)(s+j-1)-l(l+1)=(l+j+1)(l+j)-l(l+1)=j(2l+j+1)>0\quad(j\ge1).
\]
In detail, the coefficient of $r^{s+j-2}$ also receives
$\kappa Ea_{j-2}$ from $\kappa Ey$ and, since
$Vy=\sum_q\sum_ib_qa_ir^{s+i+q}$, the terms $-\kappa b_qa_i$ with
$i+q=j-2$ from $-\kappa Vy$. With negative-index coefficients zero,
this gives
\[
j(2l+j+1)a_j
=\kappa\left[\sum_{q=-1}^{j-2}b_q a_{j-2-q}-Ea_{j-2}\right].
\]
Thus an arbitrary nonzero $a_0$ determines every following coefficient;
for example $a_1=\kappa b_{-1}a_0/[2(l+1)]$. The regular solution is
$\boxed{y_1=a_0r^{l+1}[1+O(r)]}$.

(b) There is no first derivative ($P=0$, $e^{-\int P}=1$), so
reduction of order is $y_2=y_1\int^rds/y_1(s)^2$. Its leading terms are
\[
\frac1{y_1(s)^2}=a_0^{-2}s^{-2l-2}[1+O(s)],
\quad
\int^r\frac{ds}{y_1(s)^2}
\sim-\frac{a_0^{-2}}{2l+1}r^{-2l-1}:
\]
integrating termwise, $s^{-2l-2}$ gives $r^{-2l-1}/(-2l-1)$, while the
higher terms give the higher powers $r^{-2l},r^{-2l+1},\ldots$ and
possibly $\ln r$ (from an $s^{-1}$ term), all smaller as $r\to0$.
Multiplying by $y_1=a_0r^{l+1}[1+O(r)]$ gives
\[
\boxed{y_2(r)\sim-\frac1{a_0(2l+1)}r^{-l}}.
\]
\erratum{The textbook states that $y_2$ diverges at the origin as
$r^{-l}$. For $l\ge1$ it does; for $l=0$, however,
$y_2$ tends to a nonzero constant and does not diverge. The physical
radial function $R=y/r$ still behaves as $1/r$ and is irregular.}
For $l=0$ the next terms may contain $r\ln r$, which still tend to zero
and do not change the constant leading behavior.

\Needspace{5\baselineskip}
\par\bigskip
\noindent\textbf{7.6.25}\quad Show that if a second solution, $y_{2}$, is
assumed to be related to the first solution, $y_{1}$, according to
$y_{2}(x)=y_{1}(x)f(x)$, substitution back into the original equation
\[
y''_{2}+P(x)y'_{2}+Q(x)y_{2}=0
\]
leads to
\[
f(x)=\int^{x}\frac{\exp\left[-\int^{s}P(t)\,dt\right]}{[y_{1}(s)]^{2}}\,ds,
\]
in agreement with Eq.~(7.67).
\erratum{Apply the integral formula on an interval with continuous $P,Q$
and a nowhere-zero first solution $y_1$; choose all integration limits
within that interval.}
\solution
Set $y_2=y_1f$. Two differentiations give
\[
y_2'=y_1'f+y_1f',\qquad
y_2''=y_1''f+2y_1'f'+y_1f''.
\]
Substitution and use of the equation for $y_1$ leave
\[
0=f(y_1''+Py_1'+Qy_1)+(2y_1'+Py_1)f'+y_1f''
=y_1f''+(2y_1'+Py_1)f'.
\]
On an interval with $y_1\ne0$, let $g=f'$; then
\[
g'+\left(2\frac{y_1'}{y_1}+P\right)g=0,
\qquad
\frac{g'}g=-2\frac{y_1'}{y_1}-P.
\]
Integrating gives $\ln|g|=-2\ln|y_1|-\int^xP\,dt+c$ on an interval
where $g\ne0$; exponentiating, with the constant sign of $g$ absorbed,
gives $g=C e^{-\int Pdx}/y_1^2$, and a second integration gives
\[
\boxed{f=D+C\int^x\frac{e^{-\int^sP(t)dt}}{y_1(s)^2}ds}.
\]
Here $C=0$ recovers the constant-$f$ branch excluded by division by $g$.
For an independent second solution choose $C\ne0$, and omit the
$D y_1$ contribution. Taking $C=1,D=0$ gives exactly the formula in
the question; its Wronskian is
$y_1(y_1'f+y_1f')-y_1'y_1f=y_1^2f'=e^{-\int Pdx}\ne0$.

\Needspace{5\baselineskip}
\par\bigskip
\noindent\textbf{7.6.26}\quad
\erratum{Use $x>0$ and $a_0\ne0$; the primed labels in parts (b) and (c)
name limiting solutions, not derivatives.}
\begin{enumerate}
\item[(a)] Show that
  \[
  y''+\frac{1-\alpha^{2}}{4x^{2}}y=0
  \]
  has two solutions:
  \[
  y_{1}(x)=a_{0}x^{(1+\alpha)/2},\quad y_{2}(x)=a_{0}x^{(1-\alpha)/2}.
  \]
\item[(b)] For $\alpha=0$ the two linearly independent solutions of part (a)
  reduce to the single solution $y'_{1}=a_{0}x^{1/2}$. Using Eq.~(7.68) derive
  a second solution,
  \[
  y'_{2}(x)=a_{0}x^{1/2}\ln x.
  \]
  Verify that $y'_{2}$ is indeed a solution.
\item[(c)] Show that the second solution from part (b) may be obtained as a
  limiting case from the two solutions of part (a):
  \[
  y'_{2}(x)=\lim_{\alpha\to0}\left(\frac{y_{1}-y_{2}}{\alpha}\right).
  \]
\end{enumerate}
\solution
Take $x>0$ and $a_0\ne0$. (a) For a trial $y=a_0x^s$,
\[
y''+\frac{1-\alpha^2}{4x^2}y
=a_0x^{s-2}\left[s(s-1)+\frac{1-\alpha^2}{4}\right].
\]
The quadratic is
$s^2-s+\frac14-\frac{\alpha^2}4=(s-\frac12)^2-(\frac\alpha2)^2$, which
factors as
$[s-(1+\alpha)/2][s-(1-\alpha)/2]$, so both stated powers solve
the equation. Their Wronskian is
\[
W=a_0^2\left(\frac{1-\alpha}{2}-\frac{1+\alpha}{2}\right)
 x^{(1+\alpha)/2+(1-\alpha)/2-1}=-\alpha a_0^2.
\]
They are independent precisely when $\alpha\ne0$.

(b) At $\alpha=0$ the common solution is $u_1=a_0\sqrt x$;
here $u_1,u_2$ denote the primed labels in the question, not derivatives.
Since $P=0$,
\[
u_1\int^x\frac{ds}{u_1(s)^2}
=a_0\sqrt x\int^x\frac{ds}{a_0^2s}
=\frac1{a_0}\sqrt x\ln x.
\]
An arbitrary nonzero rescaling (by $a_0^2$) gives
$\boxed{u_2=a_0\sqrt x\ln x}$. Explicitly, by the product rule,
\[
u_2'=a_0\left(\frac{\ln x}{2\sqrt x}+\frac{\sqrt x}x\right)
=a_0x^{-1/2}(1+\tfrac12\ln x),
\]
\[
u_2''=a_0\left[-\tfrac12x^{-3/2}(1+\tfrac12\ln x)
+x^{-1/2}\frac1{2x}\right]
=-\frac{a_0}{4}x^{-3/2}\ln x,
\]
so
\[
u_2''+\frac{u_2}{4x^2}=-\frac{a_0}4x^{-3/2}\ln x+\frac{a_0}4x^{-3/2}\ln x=0.
\]
Also $W(u_1,u_2)=u_1u_2'-u_1'u_2
=a_0^2(1+\tfrac12\ln x)-\tfrac12a_0^2\ln x=a_0^2\ne0$.

(c) Write the powers as exponentials, using
$x^{(1\pm\alpha)/2}=\sqrt x\,e^{\pm(\alpha/2)\ln x}$:
\[
\frac{y_1-y_2}{\alpha}
=a_0\sqrt x\frac{e^{(\alpha/2)\ln x}-e^{-(\alpha/2)\ln x}}\alpha.
\]
Since $e^{hL}=1+hL+\tfrac12h^2L^2+O(h^3)$, applied with $L=\ln x$ and
$h=\pm\alpha/2$, the constant and the equal quadratic terms cancel and
the numerator is $\alpha\ln x+O(\alpha^3)$ (the odd terms double), and
hence
\[
\boxed{\lim_{\alpha\to0}\frac{y_1-y_2}{\alpha}
=a_0\sqrt x\ln x=u_2}.
\]
\erratum{The Instructor's Manual drops $a_0$ and gives an extra factor
$1/2$ in this limit. The explicit expansion above retains $a_0$ and
shows that no extra $1/2$ is present.}

\Needspace{8\baselineskip}
\subsection*{7.7 Inhomogeneous Linear ODEs}

\Needspace{5\baselineskip}
\par\bigskip
\noindent\textbf{7.7.1}\quad If our linear, second-order ODE is inhomogeneous,
that is, of the form of Eq.~(7.94), the most general solution is
\[
y(x)=\erratum{C_1}y_{1}(x)+\erratum{C_2}y_{2}(x)+y_{p}(x),
\]
where $y_{1}$ and $y_{2}$ are independent solutions of the homogeneous
equation, \erratum{on a connected interval where the normalized
coefficients and source are continuous}. Show that
\[
y_{p}(x)=y_{2}(x)\int^{x}\frac{y_{1}(s)F(s)\,ds}{W\{y_{1}(s),y_{2}(s)\}}
-y_{1}(x)\int^{x}\frac{y_{2}(s)F(s)\,ds}{W\{y_{1}(s),y_{2}(s)\}},
\]
with $W\{y_{1}(x),y_{2}(x)\}$ the Wronskian of $y_{1}(s)$ and $y_{2}(s)$.

\solution
\erratum{For fixed independent homogeneous solutions, the general solution
is $y=C_1y_1+C_2y_2+y_p$; the two arbitrary coefficients must be explicit.}
Write the normalized equation as $L[y]=y''+Py'+Qy=F$ and seek
$y_p=u_1y_1+u_2y_2$. Then
$y_p'=u_1y_1'+u_2y_2'+(u_1'y_1+u_2'y_2)$; impose the auxiliary condition
$u_1'y_1+u_2'y_2=0$, so that
\[
y_p'=u_1y_1'+u_2y_2',\qquad
y_p''=u_1y_1''+u_2y_2''+u_1'y_1'+u_2'y_2'.
\]
Collecting the terms multiplied by $u_1$ and by $u_2$,
\[
L[y_p]=u_1L[y_1]+u_2L[y_2]+u_1'y_1'+u_2'y_2'.
\]
Because $L[y_1]=L[y_2]=0$, the remaining equation is
$u_1'y_1'+u_2'y_2'=F$. Thus
\[
\begin{pmatrix}y_1&y_2\\y_1'&y_2'\end{pmatrix}
\begin{pmatrix}u_1'\\u_2'\end{pmatrix}
=\begin{pmatrix}0\\F\end{pmatrix}.
\]
Its determinant is $W=y_1y_2'-y_1'y_2$, nonzero on the regular interval:
if $W(x_0)=0$, a null vector $(a,b)\ne(0,0)$ of this matrix at $x_0$
would make $ay_1+by_2$ a homogeneous solution with zero value and slope
at $x_0$, hence identically zero by initial-value uniqueness (as in
Exercise 7.6.5), contradicting independence. Cramer's rule gives
\[
u_1'=\frac1W\det\begin{pmatrix}0&y_2\\F&y_2'\end{pmatrix}
=-\frac{y_2F}{W},\qquad
u_2'=\frac1W\det\begin{pmatrix}y_1&0\\y_1'&F\end{pmatrix}
=\frac{y_1F}{W}.
\]
Integrate and substitute:
\[
\boxed{y_p=y_2(x)\int^x\frac{y_1(s)F(s)}{W(s)}ds
-y_1(x)\int^x\frac{y_2(s)F(s)}{W(s)}ds}.
\]
Constants in the integrals contribute only multiples of $y_1,y_2$.
As a direct check, these $u_1',u_2'$ satisfy both conditions:
$u_1'y_1+u_2'y_2=(-y_2y_1+y_1y_2)F/W=0$ and
$L[y_p]=u_1'y_1'+u_2'y_2'=(y_1y_2'-y_1'y_2)F/W=F$.
Finally any other solution minus this particular solution satisfies the
homogeneous equation, and every homogeneous solution $y_h$ is a
combination $C_1y_1+C_2y_2$: by Exercise 7.6.10, $y_h,y_1,y_2$ satisfy
a nontrivial relation $cy_h+ay_1+by_2=0$, and $c\ne0$ because $y_1,y_2$
are independent. This proves that $C_1y_1+C_2y_2+y_p$ is the general
solution.

\Needspace{8\baselineskip}
\bigskip
Find the general solutions to the following inhomogeneous ODEs:

\Needspace{5\baselineskip}
\par\bigskip
\noindent\textbf{7.7.2}\quad $y''+y=1$.
\solution
The homogeneous characteristic equation is $r^2+1=0$, with roots
$r=\pm i$, giving $y_h=C_1\cos x+C_2\sin x$. For a constant particular
solution $y_p=A$, the equation requires $0+A=1$, so $A=1$. Therefore
\[
\boxed{y=C_1\cos x+C_2\sin x+1}.
\]
For a variation-of-parameters check, take $y_1=\cos x$, $y_2=\sin x$,
$W=\cos x\cos x-(-\sin x)\sin x=1$, $F=1$. Then, by Exercise 7.7.1,
$u_1'=-y_2F/W=-\sin x$, $u_2'=y_1F/W=\cos x$, so one choice is
$u_1=\cos x$, $u_2=\sin x$. Their combination is
$y_p=u_1y_1+u_2y_2=\cos^2x+\sin^2x=1$, agreeing with the constant trial.
\erratum{The Instructor's Manual labels both derivatives $u_1'$ in its
variation-of-parameters calculation; the second must be
$u_2'=\cos x$, as Cramer's rule above shows.}

\Needspace{5\baselineskip}
\par\bigskip
\noindent\textbf{7.7.3}\quad $y''+4y=e^{x}$.
\solution
The homogeneous roots satisfy $r^2+4=0$, so $r=\pm2i$ and the real
homogeneous solution is $C_1\cos2x+C_2\sin2x$. Try $y_p=Ae^x$; then
\[
y_p''+4y_p=(A+4A)e^x=5Ae^x.
\]
Matching the right side sets $A=1/5$. Thus
\[
\boxed{y=C_1\cos2x+C_2\sin2x+\frac15e^x}.
\]
The homogeneous Wronskian is
$W(\cos2x,\sin2x)=2\cos^22x+2\sin^22x=2$, so the two constants are
independent. Differentiating the displayed solution twice verifies that
its trigonometric contributions cancel and its exponential contribution
is exactly $e^x$.

Variation of parameters (Exercise 7.7.1) gives the same particular
solution: with $y_1=\cos2x$, $y_2=\sin2x$, $W=2$, $F=e^x$,
\[
u_1'=-\frac{y_2F}W=-\frac{e^x\sin2x}2,\qquad
u_2'=\frac{y_1F}W=\frac{e^x\cos2x}2.
\]
The real and imaginary parts of
$\int e^{(1+2i)x}dx=e^{(1+2i)x}/(1+2i)=e^x(\cos2x+i\sin2x)(1-2i)/5$
give $\int e^x\cos2x\,dx=e^x(\cos2x+2\sin2x)/5$ and
$\int e^x\sin2x\,dx=e^x(\sin2x-2\cos2x)/5$, so
\[
u_1=\frac{e^x(2\cos2x-\sin2x)}{10},\qquad
u_2=\frac{e^x(\cos2x+2\sin2x)}{10},
\]
and $y_p=u_1\cos2x+u_2\sin2x=e^x(2\cos^22x+2\sin^22x)/10=e^x/5$.
\erratum{In the Instructor's Manual's variation-of-parameters
calculation, $u_2'=e^x\cos(2x)/2$, not $e^x\sin(2x)/2$: Cramer's rule
gives $y_1F/W=\cos(2x)e^x/2$. Its printed antiderivative and final
answer agree with this corrected derivative.}

\Needspace{5\baselineskip}
\par\bigskip
\noindent\textbf{7.7.4}\quad $y''-3y'+2y=\sin x$.
\solution
The homogeneous characteristic polynomial is
$r^2-3r+2=(r-1)(r-2)$, giving $y_h=C_1e^x+C_2e^{2x}$.
For $y_p=A\sin x+B\cos x$,
\[
y_p'=A\cos x-B\sin x,\qquad y_p''=-A\sin x-B\cos x.
\]
Hence, collecting $\sin x$ terms ($-A+3B+2A$) and $\cos x$ terms
($-B-3A+2B$),
\[
y_p''-3y_p'+2y_p=(A+3B)\sin x+(B-3A)\cos x.
\]
Matching coefficients gives $A+3B=1$, $B-3A=0$; hence
$B=3A$, $10A=1$, $A=1/10$, $B=3/10$. The general solution is
\[
\boxed{y=C_1e^x+C_2e^{2x}+\frac{\sin x+3\cos x}{10}}.
\]
The coefficient equations above are also a direct residual check of the
particular solution; the homogeneous Wronskian is $e^{3x}\ne0$.

\Needspace{5\baselineskip}
\par\bigskip
\noindent\textbf{7.7.5}\quad $xy''-(1+x)y'+y=x^{2}$.
\solution
First find homogeneous solutions. For a linear trial $y=ax+b$,
$xy''-(1+x)y'+y=b-a$, so $y_1=x+1$ is a solution.
The exponential $y_2=e^x$ also works, since
$xe^x-(1+x)e^x+e^x=0$. Their Wronskian is
\[
W=(x+1)e^x-e^x=xe^x,
\]
nonzero on either interval $x>0$ or $x<0$.
After division by $x$ the equation reads
$y''-\frac{1+x}xy'+\frac1xy=x$, so the source term for variation of
parameters is $F=x$, not $x^2$. Therefore
\[
u_1'=-\frac{y_2F}W=-\frac{e^x x}{xe^x}=-1,\qquad
u_2'=\frac{y_1F}W=\frac{(x+1)x}{xe^x}=(x+1)e^{-x}.
\]
Integration by parts gives
$\int(x+1)e^{-x}dx=-(x+1)e^{-x}+\int e^{-x}dx=-(x+2)e^{-x}$, so one
choice of antiderivatives is
$u_1=-x$, $u_2=-(x+2)e^{-x}$, giving
\[
y_p=u_1y_1+u_2y_2=-x(x+1)-(x+2)=-x^2-2x-2.
\]
Since $-2(x+1)$ is homogeneous, it can be absorbed into its arbitrary
coefficient; take simply $y_p=-x^2$. Thus
\[
\boxed{y=C_1(x+1)+C_2e^x-x^2}.
\]
The direct check for the particular solution, with $y_p'=-2x$ and
$y_p''=-2$, is $x(-2)-(1+x)(-2x)-x^2=-2x+2x+2x^2-x^2=x^2$.
The formula extends smoothly through zero, although the normalized ODE
and the Wronskian argument are applied on intervals away from zero.
It also describes every $C^2$ solution across zero: write the formula
with constants $C_1^\pm,C_2^\pm$ on the two sides. Continuity of $y''$
gives $C_2^-=C_2^+$ because $y''(0)=C_2-2$; continuity of $y$ then
gives $C_1^-=C_1^+$. Thus no additional $C^2$ branch is lost at zero.

\Needspace{8\baselineskip}
\subsection*{7.8 Nonlinear Differential Equations}

\Needspace{5\baselineskip}
\par\bigskip
\noindent\textbf{7.8.1}\quad Consider the Riccati equation
$y'=y^{2}-y-2$. A particular solution to this equation is $y=2$. Find a more
general solution.
\solution
Use the given solution by setting $y=2+u$. Then
\[
u'=(2+u)^2-(2+u)-2=u^2+3u.
\]
The branch $u=0$ gives the original solution $y=2$. Since $u^2+3u$ is
smooth, initial-value problems for $u$ have unique solutions, so a
solution with $u=0$ at one point is this zero solution; any other
solution has $u\ne0$ throughout its interval. For such $u$ put
$v=1/u$; then
\[
v'=-\frac{u'}{u^2}=-\frac{u^2+3u}{u^2}=-1-3v,\qquad
(e^{3x}v)'=e^{3x}(v'+3v)=-e^{3x},
\]
the second form after multiplying by the integrating factor $e^{3x}$.
Integrating gives $e^{3x}v=-\frac13e^{3x}+C$, so $v=Ce^{-3x}-1/3$ and
$u=1/v=3/(3Ce^{-3x}-1)$. Renaming $3C$ as $K$,
\[
\boxed{y=2+\frac3{Ke^{-3x}-1}},\qquad \boxed{y\equiv2}.
\]
The constant solution $y=-1$ is already included at $K=0$.
Every nonconstant expression is used on a maximal interval where its
denominator is nonzero. To check it, write $D=Ke^{-3x}-1$, so $u=3/D$
and $D'=-3Ke^{-3x}$:
$u'=-3D'/D^2=9Ke^{-3x}/D^2$, while
$u^2+3u=(9+9D)/D^2=9(1+D)/D^2=9Ke^{-3x}/D^2$.
\erratum{The integrating-constant formula for $v$ must have the particular
term $-1/3$; the Instructor's Manual's intermediate $+1/3$ does not
satisfy $v'+3v=-1$.}

\Needspace{5\baselineskip}
\par\bigskip
\noindent\textbf{7.8.2}\quad A particular solution to
$y'=y^{2}/x^{3}-y/x+2x$ is $y=x^{2}$. Find a more general solution.
\solution
Work on $x\ne0$ and set $y=x^2+u$. Expanding every term gives
\[
2x+u'=\frac{x^4+2x^2u+u^2}{x^3}
          -\frac{x^2+u}{x}+2x
=2x+\frac ux+\frac{u^2}{x^3}.
\]
Thus $u'=u/x+u^2/x^3$. The branch $u=0$ is $y=x^2$. Since the right
side is smooth in $u$ for $x\ne0$, uniqueness of initial-value problems
(as in Exercise 7.8.1) means that any other solution has $u\ne0$
throughout its interval. For such $u$ let $v=1/u$; then
$v'=-u'/u^2=-1/(xu)-1/x^3$, that is,
\[
v'=-\frac vx-\frac1{x^3},\qquad
v'+\frac vx=-\frac1{x^3},\qquad (xv)'=-\frac1{x^2},
\]
the last form after multiplying by the integrating factor
$e^{\int dx/x}=x$ (up to sign).
Integrating gives $xv=1/x+C$, so $v=(1+Cx)/x^2$ and
$u=1/v=x^2/(1+Cx)$. Therefore
\[
\boxed{y=x^2+\frac{x^2}{1+Cx}},\qquad \boxed{y=x^2}.
\]
The first family is valid where $x\ne0$ and $1+Cx\ne0$; the second is
the lost $u=0$ branch. For a check,
$u=x^2/(1+Cx)$ gives, by the quotient rule,
\[
u'=\frac{2x(1+Cx)-Cx^2}{(1+Cx)^2}=\frac{2x+Cx^2}{(1+Cx)^2},\qquad
\frac ux+\frac{u^2}{x^3}=\frac x{1+Cx}+\frac x{(1+Cx)^2}
=\frac{2x+Cx^2}{(1+Cx)^2}.
\]
\erratum{The Instructor's Manual's intermediate linear equation is
missing the factor $1/x$ multiplying $v$; its displayed final solution
agrees with the corrected equation derived here.}

\Needspace{5\baselineskip}
\par\bigskip
\noindent\textbf{7.8.3}\quad Solve the Bernoulli equation $y'+xy=xy^{3}$.
\solution
The solution $y=0$ must be recorded before dividing by a power of $y$.
Since $xy^3-xy$ is smooth in $y$, initial-value problems for this
equation have unique solutions, so a solution that vanishes at one
point is the zero solution. Every other solution is nowhere zero and,
being continuous, keeps one sign on its interval. For such a solution
set $u=y^{-2}$. Then
\[
u'=-2y^{-3}y'=-2y^{-3}(xy^3-xy)=-2x+2xu.
\]
Thus, multiplying the linear equation by the integrating factor
$e^{-x^2}$,
\[
u'-2xu=-2x,\qquad (e^{-x^2}u)'=-2xe^{-x^2}=(e^{-x^2})'.
\]
After integration $e^{-x^2}u=e^{-x^2}+C$, so $u=1+Ce^{x^2}$.
The real solutions are
\[
\boxed{y=0\quad\hbox{or}\quad
 y=\pm\frac1{\sqrt{1+Ce^{x^2}}}}.
\]
\erratum{Both signs and the zero solution are needed; taking only the
positive square root loses real solutions.} Choose a connected interval
where $1+Ce^{x^2}>0$ and a fixed sign. In particular $C=0$ gives
$y=\pm1$. Differentiation of $y=\pm(1+Ce^{x^2})^{-1/2}$ gives the
first line below, and the right side of the equation, written as
$y'=xy^3-xy$, gives the same expression:
\[
\begin{aligned}
y'&=\mp\tfrac12(1+Ce^{x^2})^{-3/2}\,2xCe^{x^2}
=-\frac{xCe^{x^2}y}{1+Ce^{x^2}},\\
xy^3-xy&=xy\left[\frac1{1+Ce^{x^2}}-1\right]
=-\frac{xCe^{x^2}y}{1+Ce^{x^2}}.
\end{aligned}
\]

\Needspace{5\baselineskip}
\par\bigskip
\noindent\textbf{7.8.4}\quad ODEs of the form $y=xy'+f(y')$ are known as
Clairaut equations. The first step in solving an equation of this type is to
differentiate it, \erratum{assuming a twice differentiable solution}, yielding
\[
y'=y'+xy''+f'(y')y'',\quad\text{or }y''\left[x+f'(y')\right]=0.
\]
Solutions may therefore be obtained both from $y''=0$ and from
$f'(y')=-x$. The so-called general solution comes from $y''=0$.

For $f(y')=(y')^{2}$,
\begin{enumerate}
\item[(a)] Obtain the general solution (note that it contains a single
  constant).
\item[(b)] Obtain the so-called singular solution from $f'(y')=-x$. By
  substituting back into the original ODE show that this singular solution
  contains no adjustable constants.
\end{enumerate}
\emph{Note.}~The singular solution is the envelope of the general solutions.
\solution
For $f(p)=p^2$, differentiating the equation gives
$y''(x+2y')=0$.

(a) On the branch $y''=0$, write $y=ax+b$ and $y'=a$.
Substitute into the original, rather than only its derivative:
\[
ax+b=xa+a^2\quad\Longrightarrow\quad b=a^2.
\]
Hence the one-parameter family is
\[
\boxed{y=ax+a^2},\qquad a\text{ arbitrary}.
\]
(b) On the other branch $x+2y'=0$, so $y'=-x/2$ and
$y=-x^2/4+C$. The original equation requires
\[
-\frac{x^2}4+C=x\left(-\frac x2\right)+\left(-\frac x2\right)^2
=-\frac{x^2}4,
\]
so $C=0$ and
\[
\boxed{y=-x^2/4}.
\]
To verify the envelope statement, eliminate the parameter $a$ from
$y=ax+a^2$ and its parameter derivative $0=x+2a$.
This gives $a=-x/2$ and $y=-x^2/4$. Each line touches the parabola
at $x=-2a$: both values there are $-a^2$ and both slopes are $a$.

\erratum{If only $C^1$ solutions of the original first-order equation
are required, the two types can also be joined at a contact point.}
For example, $y=0$ for $x\le0$ and $y=-x^2/4$ for $x\ge0$ has
continuous value and derivative at zero, and each branch satisfies
$y=xy'+(y')^2$. The second derivatives there are $0$ and $-1/2$,
so this joined solution lies outside the differentiation argument.
The requested line family and singular curve are the standard smooth
branches; they do not exhaust such $C^1$ joins.
